\documentclass{article}
\usepackage[margin=0.8in]{geometry}
\usepackage{amsmath,amssymb,amsthm,mathtools}
\usepackage{mathrsfs,bm,bbm}
\usepackage{graphicx,booktabs,array,multirow,tabularx,longtable}
\usepackage{enumitem}
\usepackage{xcolor}
\usepackage{algorithm,algpseudocode}
\usepackage{subcaption,tikz}
\usepackage{siunitx,cancel,accents}
\usepackage{microtype}
\usepackage[numbers,sort&compress]{natbib}
\usepackage[colorlinks=true,linkcolor=blue,citecolor=blue,urlcolor=blue]{hyperref}
\usepackage{cleveref}
\setlength{\emergencystretch}{2em}
\newtheorem{assumption}{Assumption}
\newtheorem{definition}{Definition}
\newtheorem{theorem}{Theorem}
\newtheorem{lemma}{Lemma}
\newtheorem{proposition}{Proposition}
\newtheorem{openendedquestion}{Open-ended Question}
\newtheorem{conjecture}{Conjecture}
\newcommand{\coreref}[1]{\ref{#1}}

\begin{document}

\sloppy

\section*{exp\_\allowbreak{}randomregular\_\allowbreak{}fullgate\_\allowbreak{}stableset\_\allowbreak{}frontier}

\noindent\textit{Specialization:} v1\par

\paragraph{TL;DR.} For fixed potential outcomes under arbitrary neighborhood interference on a revealed random \(d\)-regular graph, the fully audited unconditional comparison is the published-order bracket: uniformly for feasible \(d=o(\sqrt n)\), with probability \(1-o(1)\), \(c(d+1)/n\le R_2(G)\le C(1+d^2)/n\). For fixed \(d\), this gives \(R_2(G)=\Theta_d(n^{-1})\). The proposed stronger lower edge \(r(n,d)=(d+1)^2/[n\log(d+2)]\) still requires a valid random-regular balanced-hole theorem. The reservoir capsule has exact covariance identities and an exact finite audit conditional on positive focal inclusion, but its \(O(n/k)\) norm bound, scalable compiler, and logarithmically improved risk remain open.

\section{Project justification}

\paragraph{Gap.} Published unrestricted minimax theory places growing-degree random-regular full-GATE risk between orders \(d/n\) and \(d^2/n\). A deterministic balanced-hole Bayes reduction can identify the graph certificate needed for the stronger lower scale, but the required uniform random-regular switching comparison has not been proved or anchored by a matching primary theorem.

\paragraph{Niche.} The two-copy exposure graph and balanced reservoir capsule expose exact covariance structure and a concrete computational target. They isolate two unresolved tasks: the accepted compatible-pair operator bound and a polynomial-time observed-data compiler with auditable error.

\paragraph{Fill.} The paper retains the supported all-procedure lower comparison and the observed-data adaptation of the published conflict-graph upper design, records the \(\Theta_d(n^{-1})\) fixed-degree benchmark, proves the capsule's exact finite covariance algebra, and supplies a finite evaluator that detects zero inclusion before computing kernel or estimator formulas. It does not claim an exact frontier or capsule attainment.

\section{Related work}

Kandiros, Harshaw, and Savje define the same endpoint-\(L^2\) unrestricted minimax experiment and prove the growing-degree random-regular bracket \(d/n\lesssim R_2(G)\lesssim d^2/n\) under \(d=o(\sqrt n)\). The attempted strengthening of its lower edge to \(r(n,d)\) reduces to a balanced-hole estimate, but the needed simple-random-regular switching counts are not yet verified. Kandiros, Pipis, Daskalakis, and Harshaw provide the closed upper route used here: Lemma 3.3 gives exact desired-exposure probabilities, Lemma 4.2 controls local normalized covariances, and Theorem 4.1 gives the \(O(\lambda(H_G)/n)\) modified Horvitz--Thompson variance bound. Li and Li prove only a treated-component result. Harshaw, S\"avje, and Wang motivate operator-norm risk analysis, while Perkins and Wang and Cai, Zhang, and Airoldi study balanced or independent sets in different graph and estimand settings; none supplies the acceptance-conditioned two-arm capsule covariance estimate or compiler required here.

\section{Setup}

Target (estimand / structural parameter): \(\tau_G(y)=n^{-1}\sum_{i\in V}\{y_i(\boldsymbol 1)-y_i(\boldsymbol 0)\}\).
Primary functional / estimator / diagnostic: Uniformly for feasible \(d=o(\sqrt n)\), with probability \(1-o(1)\), \(c(1+d)/n\le R_2(G)\le C(1+d^2)/n\); for every fixed feasible \(d\), \(R_2(G)=\Theta_d(n^{-1})\). The upper bound is attained by the published Conflict Graph Design after observed-data Rao--Blackwellization, while the stronger balanced-hole lower edge and logarithmic capsule rate remain open..
\begin{itemize}
  \item \texttt{n} --- \texttt{population size}; \(\{4,5,\ldots\}\); \texttt{asymptotic index}
  \par\smallskip\noindent\emph{Definition.} \texttt{The triangular-\allowbreak{}array index and number of experimental units.}
  \item \texttt{d} --- \texttt{regular degree}; \(\{0,\ldots,n-1\}\); \texttt{structural regime parameter}
  \par\smallskip\noindent\emph{Definition.} The degree at index \(n\), with \(nd\) even.
  \item \texttt{V} --- \texttt{finite population}; \(\{1,\ldots,n\}\); \texttt{population}
  \par\smallskip\noindent\emph{Definition.} \texttt{The labeled experimental units.}
  \item \texttt{i} --- \texttt{unit index}; \(V\); \texttt{index}
  \par\smallskip\noindent\emph{Definition.} \texttt{A generic experimental unit.}
  \item \texttt{j} --- \texttt{unit index}; \(V\); \texttt{index}
  \par\smallskip\noindent\emph{Definition.} \texttt{A second generic experimental unit.}
  \item \texttt{a} --- \texttt{exposure-\allowbreak{}arm index}; \(\{0,1\}\); \texttt{index}
  \par\smallskip\noindent\emph{Definition.} \texttt{A generic endpoint arm.}
  \item \texttt{b} --- \texttt{exposure-\allowbreak{}arm index}; \(\{0,1\}\); \texttt{index}
  \par\smallskip\noindent\emph{Definition.} \texttt{A second endpoint arm.}
  \item \texttt{G\_\allowbreak{}n} --- \texttt{random interference graph}; simple labeled graphs on \(V\); \texttt{random environment}
  \par\smallskip\noindent\emph{Definition.} A uniformly random simple labeled \(d\)-regular graph on \(V\).
  \item \texttt{G} --- \texttt{revealed interference graph}; simple labeled \(d\)-regular graphs on \(V\); \texttt{design input}
  \par\smallskip\noindent\emph{Definition.} A realization of \(G_n\), observed before the design is drawn.
  \item \texttt{\textbackslash{}\allowbreak{}operatorname\{dist\}\allowbreak{}\_\allowbreak{}G(i,\allowbreak{}j)} --- \texttt{graph distance}; \(\mathbb N\cup\{\infty\}\); \texttt{conflict geometry}
  \par\smallskip\noindent\emph{Definition.} The shortest-path distance between \(i\) and \(j\) in \(G\).
  \item \texttt{N\_\allowbreak{}G[i]} --- \texttt{closed neighborhood}; \(2^V\); \texttt{interference support}
  \par\smallskip\noindent\emph{Definition.} The set \(N_G[i]=\{j\in V:\operatorname{dist}_G(i,j)\le 1\}\).
  \item \texttt{z} --- \texttt{assignment argument}; \(\{0,1\}^V\); \texttt{potential-\allowbreak{}outcome argument}
  \par\smallskip\noindent\emph{Definition.} \texttt{A generic binary treatment assignment.}
  \item \texttt{z'} --- \texttt{assignment argument}; \(\{0,1\}^V\); \texttt{potential-\allowbreak{}outcome argument}
  \par\smallskip\noindent\emph{Definition.} \texttt{A second binary treatment assignment.}
  \item \texttt{Z} --- \texttt{random assignment}; \(\{0,1\}^V\); \texttt{observed assignment}
  \par\smallskip\noindent\emph{Definition.} \texttt{The random assignment generated by a generic experimental design.}
  \item \texttt{\textbackslash{}\allowbreak{}boldsymbol 0} --- \texttt{all-\allowbreak{}control assignment}; \(\{0,1\}^V\); \texttt{endpoint assignment}
  \par\smallskip\noindent\emph{Definition.} \texttt{The assignment sending every unit to control.}
  \item \texttt{\textbackslash{}\allowbreak{}boldsymbol 1} --- \texttt{all-\allowbreak{}treated assignment}; \(\{0,1\}^V\); \texttt{endpoint assignment}
  \par\smallskip\noindent\emph{Definition.} \texttt{The assignment sending every unit to treatment.}
  \item \texttt{y} --- \texttt{potential-\allowbreak{}outcome schedule}; \((\mathbb R^V)^{\{0,1\}^V}\); \texttt{adversarial schedule}
  \par\smallskip\noindent\emph{Definition.} \texttt{A fixed finite-\allowbreak{}population potential-\allowbreak{}outcome schedule.}
  \item \texttt{y\_\allowbreak{}i(z)} --- \texttt{potential outcome}; \(\mathbb R\); \texttt{causal primitive}
  \par\smallskip\noindent\emph{Definition.} The outcome of unit \(i\) under assignment \(z\).
  \item \texttt{y\_\allowbreak{}\{i1\}\allowbreak{}} --- \texttt{treated-\allowbreak{}endpoint outcome}; \(\mathbb R\); \texttt{estimand component}
  \par\smallskip\noindent\emph{Definition.} The endpoint value \(y_{i1}=y_i(\boldsymbol 1)\).
  \item \texttt{y\_\allowbreak{}\{i0\}\allowbreak{}} --- \texttt{control-\allowbreak{}endpoint outcome}; \(\mathbb R\); \texttt{estimand component}
  \par\smallskip\noindent\emph{Definition.} The endpoint value \(y_{i0}=y_i(\boldsymbol 0)\).
  \item \texttt{y\_\allowbreak{}\{ia\}\allowbreak{}} --- \texttt{arm-\allowbreak{}indexed endpoint outcome}; \(\mathbb R\); \texttt{contrast component}
  \par\smallskip\noindent\emph{Definition.} The value \(y_{ia}=y_{i1}\) for \(a=1\) and \(y_{ia}=y_{i0}\) for \(a=0\).
  \item \texttt{\textbackslash{}\allowbreak{}tau\_\allowbreak{}G} --- \texttt{finite-\allowbreak{}population causal estimand}; \(\mathbb R\); \texttt{target}
  \par\smallskip\noindent\emph{Definition.} The global average treatment effect \(\tau_G(y)=n^{-1}\sum_{i\in V}(y_{i1}-y_{i0})\).
  \item \texttt{\textbackslash{}\allowbreak{}mathcal U} --- \texttt{endpoint-\allowbreak{}exposure vertex set}; \(V\times\{0,1\}\); \texttt{conflict-\allowbreak{}graph vertices}
  \par\smallskip\noindent\emph{Definition.} \texttt{The two labeled endpoint exposures for every unit.}
  \item \texttt{u} --- \texttt{endpoint-\allowbreak{}exposure vertex}; \(\mathcal U\); \texttt{index}
  \par\smallskip\noindent\emph{Definition.} A generic vertex \(u=(i,a)\).
  \item \texttt{v} --- \texttt{endpoint-\allowbreak{}exposure vertex}; \(\mathcal U\); \texttt{index}
  \par\smallskip\noindent\emph{Definition.} A second generic vertex \(v=(j,b)\).
  \item \texttt{s\_\allowbreak{}a} --- \texttt{contrast sign}; \(\{-1,1\}\); \texttt{contrast orientation}
  \par\smallskip\noindent\emph{Definition.} The sign \(s_1=1\) and \(s_0=-1\).
  \item \texttt{H\_\allowbreak{}G} --- \texttt{two-\allowbreak{}copy GATE conflict graph}; simple bipartite graphs on \(\mathcal U\); \texttt{exposure incompatibility}
  \par\smallskip\noindent\emph{Definition.} The graph joining \((i,a)\) and \((j,b)\) exactly when \(a\ne b\) and \(\operatorname{dist}_G(i,j)\le 2\).
  \item \texttt{\textbackslash{}\allowbreak{}lambda(H\_\allowbreak{}G)} --- \texttt{conflict-\allowbreak{}graph spectral radius}; \([0,\infty)\); \texttt{published-\allowbreak{}design comparator}
  \par\smallskip\noindent\emph{Definition.} The largest eigenvalue of the adjacency matrix of \(H_G\).
  \item \texttt{k} --- \texttt{focal arm size}; \(\{1,\ldots,\lfloor n/4\rfloor\}\); \texttt{design tuning}
  \par\smallskip\noindent\emph{Definition.} The integer \(k=\max\{1,\lfloor n\log(d+2)/[32(d+1)^2]\rfloor\}\wedge\lfloor n/4\rfloor\).
  \item \texttt{A} --- \texttt{first focal reservoir}; \(\{S\subseteq V:|S|=k\}\); \texttt{design random variable}
  \par\smallskip\noindent\emph{Definition.} A uniformly sampled \(k\)-subset before the capsule acceptance check.
  \item \texttt{F\_\allowbreak{}G(A)} --- \texttt{opposite-\allowbreak{}arm reservoir}; \(2^V\); \texttt{compatible controls}
  \par\smallskip\noindent\emph{Definition.} The set \(F_G(A)=\{j\in V:\min_{i\in A}\operatorname{dist}_G(i,j)>2\}\).
  \item \texttt{C} --- \texttt{second focal reservoir}; \(\{S\subseteq V:|S|=k\}\); \texttt{design random variable}
  \par\smallskip\noindent\emph{Definition.} Conditional on \(|F_G(A)|\ge k\), a uniformly sampled \(k\)-subset of \(F_G(A)\).
  \item \texttt{\textbackslash{}\allowbreak{}omega} --- \texttt{arm-\allowbreak{}swap coin}; \(\operatorname{Bernoulli}(1/2)\); \texttt{arm symmetry}
  \par\smallskip\noindent\emph{Definition.} \texttt{A fair coin independent of the reservoir draws.}
  \item \texttt{T} --- \texttt{treated focal set}; \(2^V\); \texttt{treated endpoint sample}
  \par\smallskip\noindent\emph{Definition.} The set \(T=A\) if \(\omega=1\) and \(T=C\) if \(\omega=0\).
  \item \texttt{C\_\allowbreak{}0} --- \texttt{control focal set}; \(2^V\); \texttt{control endpoint sample}
  \par\smallskip\noindent\emph{Definition.} The set \(C_0=C\) if \(\omega=1\) and \(C_0=A\) if \(\omega=0\).
  \item \texttt{Z\textasciicircum{}\{\textbackslash{}\allowbreak{}mathrm\{cap\}\allowbreak{}\}\allowbreak{}} --- \texttt{random assignment}; \(\{0,1\}^V\); \texttt{observed randomization}
  \par\smallskip\noindent\emph{Definition.} The assignment that treats \(\bigcup_{i\in T}N_G[i]\), controls \(\bigcup_{i\in C_0}N_G[i]\), and assigns independent fair coins to all remaining units.
  \item \texttt{\textbackslash{}\allowbreak{}mathsf D\_\allowbreak{}G\textasciicircum{}\{\textbackslash{}\allowbreak{}mathrm\{cap\}\allowbreak{}\}\allowbreak{}} --- \texttt{capsule design law}; probability measures on \(\{0,1\}^V\); \texttt{chosen design}
  \par\smallskip\noindent\emph{Definition.} The law generated by at most \(n^3\) independent trials of \(A\), accepting the first with \(|F_G(A)|\ge k\), then drawing \(C\), \(\omega\), and \(Z^{\mathrm{cap}}\); if no trial is accepted, it assigns all units according to one fair global coin.
  \item \texttt{p\_\allowbreak{}\{\textbackslash{}\allowbreak{}mathrm\{acc\}\allowbreak{}\}\allowbreak{}(G)} --- \texttt{single-\allowbreak{}trial capsule acceptance probability}; \([0,1]\); \texttt{sampler certificate}
  \par\smallskip\noindent\emph{Definition.} The conditional probability \(\Pr\{|F_G(A)|\ge k\mid G\}\) for a uniform \(A\).
  \item \texttt{X\_\allowbreak{}u} --- \texttt{focal exposure indicator}; \(\{0,1\}\); \texttt{estimator inclusion}
  \par\smallskip\noindent\emph{Definition.} On an accepted capsule, \(X_{(i,1)}=\mathbf 1\{i\in T\}\) and \(X_{(i,0)}=\mathbf 1\{i\in C_0\}\); on the capped fallback, \(X_u=0\).
  \item \texttt{\textbackslash{}\allowbreak{}pi\_\allowbreak{}u} --- \texttt{focal inclusion probability}; \([0,1]\); \texttt{inverse-\allowbreak{}probability weight}
  \par\smallskip\noindent\emph{Definition.} The probability \(\pi_u=\mathbb E_{\mathsf D_G^{\mathrm{cap}}}[X_u\mid G]\).
  \item \texttt{\textbackslash{}\allowbreak{}Gamma\_\allowbreak{}G} --- \texttt{signed normalized exposure-\allowbreak{}covariance kernel}; \(\mathbb R^{\mathcal U\times\mathcal U}\); \texttt{optimality certificate}
  \par\smallskip\noindent\emph{Definition.} For \(u=(i,a)\) and \(v=(j,b)\), \((\Gamma_G)_{uv}=s_as_b\operatorname{Cov}(X_u/\pi_u,X_v/\pi_v\mid G)\) whenever all inclusion probabilities are positive.
  \item \texttt{\textbackslash{}\allowbreak{}widetilde\textbackslash{}\allowbreak{}tau\_\allowbreak{}G} --- \texttt{raw focal Horvitz-\allowbreak{}-\allowbreak{}Thompson estimator}; \(\mathbb R\); \texttt{analysis estimator}
  \par\smallskip\noindent\emph{Definition.} The estimator \(\widetilde\tau_G=n^{-1}\sum_{i\in V}\sum_{a\in\{0,1\}}s_aX_{(i,a)}y_{ia}/\pi_{(i,a)}\) when all \(\pi_u>0\), and zero otherwise.
  \item \texttt{\textbackslash{}\allowbreak{}widehat\textbackslash{}\allowbreak{}tau\_\allowbreak{}G} --- \texttt{observed-\allowbreak{}data estimator}; \(\mathbb R\); \texttt{implementable estimator}
  \par\smallskip\noindent\emph{Definition.} The Rao--Blackwell estimator \(\widehat\tau_G=\mathbb E[\widetilde\tau_G\mid G,Z^{\mathrm{cap}},y(Z^{\mathrm{cap}})]\).
  \item \texttt{\textbackslash{}\allowbreak{}mathcal M\_\allowbreak{}2(G)} --- endpoint-\(L^2\) ANI class; \texttt{potential-\allowbreak{}outcome schedules}; \texttt{minimax outcome class}
  \par\smallskip\noindent\emph{Definition.} The class constructed in \(\text{def:endpoint-class}\).
  \item \texttt{R\_\allowbreak{}2(G)} --- \texttt{unrestricted design-\allowbreak{}based minimax risk}; \([0,\infty)\); \texttt{minimax benchmark}
  \par\smallskip\noindent\emph{Definition.} The quantity \(R_2(G)=\inf_{(\mathsf D,T_*)}\sup_{y\in\mathcal M_2(G)}\mathbb E_{\mathsf D}[(T_*(G,Z,y(Z))-\tau_G(y))^2\mid G]\), where \(\mathsf D\) ranges over outcome-independent assignment laws and \(T_*\) over all measurable possibly randomized estimators.
  \item \texttt{r(n,\allowbreak{}d)} --- \texttt{candidate minimax scale}; \((0,\infty)\); \texttt{frontier scale}
  \par\smallskip\noindent\emph{Definition.} The deterministic scale \(r(n,d)=(d+1)^2/[n\log(d+2)]\).
  \item \texttt{\textbackslash{}\allowbreak{}mathsf H\_\allowbreak{}\{\textbackslash{}\allowbreak{}mathrm\{loc\}\allowbreak{}\}\allowbreak{}} --- \texttt{local-\allowbreak{}covariance compiler handle}; \texttt{randomized graph algorithms}; \texttt{open construction handle}
  \par\smallskip\noindent\emph{Definition.} Use truncated radius-two avoidance messages and randomized probing of the capsule law to target the inclusion vector \((\pi_u)_{u\in\mathcal U}\), the action of \(\Gamma_G\), and the Rao--Blackwell weights without enumerating reservoir pairs.
  \item \texttt{U\_\allowbreak{}G} --- \texttt{graph-\allowbreak{}measurable deterministic risk certificate}; \([0,\infty)\); \texttt{open compiler output}
  \par\smallskip\noindent\emph{Definition.} A deterministic function \(U_G=U_G(G)\) of the revealed graph, to be derived by \(\text{oeq:local-compiler}\).
\end{itemize}

\section{Assumptions}

\begin{assumption}[ass:degree-diverges]\label{ass:degree-diverges}
Along the triangular array, \(d=d(n)\to\infty\).
\par\smallskip\noindent\emph{Novel.} The focused growing-degree regime isolates sparse networks whose local interference burden increases with population size; sequences such as \(d=n^\alpha\) with \(0<\alpha<1/2\) are concrete members. Divergence of \(d\) is required for the claimed factor \(\log(d+2)\) improvement to be asymptotically meaningful and is not attributed to the published random-regular theory.
\end{assumption}

\begin{assumption}[ass:subsqrt-degree]\label{ass:subsqrt-degree}
Along the triangular array, \(d/\sqrt n\to0\).
\par\smallskip\noindent\emph{Standard} (Random-regular full-GATE sparsity regime, \cite{KandirosHarshawSavje2026DesignBasedMinimax}).
\end{assumption}

\begin{assumption}[ass:ani]\label{ass:ani}
For every \(i\in V\) and \(z,z'\in\{0,1\}^V\), equality of \(z\) and \(z'\) on \(N_G[i]\) implies \(y_i(z)=y_i(z')\).
\par\smallskip\noindent\emph{Standard} (Arbitrary neighborhood interference, \cite{KandirosHarshawSavje2026DesignBasedMinimax}).
\end{assumption}

\begin{assumption}[ass:treated-endpoint-l2]\label{ass:treated-endpoint-l2}
The treated endpoints satisfy \(n^{-1}\sum_{i\in V}y_{i1}^2\le1\).
\par\smallskip\noindent\emph{Standard} (Empirical endpoint second-moment restriction, \cite{KandirosHarshawSavje2026DesignBasedMinimax}).
\end{assumption}

\begin{assumption}[ass:control-endpoint-l2]\label{ass:control-endpoint-l2}
The control endpoints satisfy \(n^{-1}\sum_{i\in V}y_{i0}^2\le1\).
\par\smallskip\noindent\emph{Standard} (Empirical endpoint second-moment restriction, \cite{KandirosHarshawSavje2026DesignBasedMinimax}).
\end{assumption}

\section{Definitions}

\begin{definition}[def:endpoint-class]\label{def:endpoint-class}
\par\smallskip\noindent\emph{Defined object.} Endpoint-\(L^2\) ANI class
\par\smallskip\noindent\emph{Construction.} \(\mathcal M_2(G)=\{y^{(n)}:\text{for every index }n,\ y_i(z)=y_i(z')\text{ whenever }z|_{N_G[i]}=z'|_{N_G[i]},\ n^{-1}\sum_{i\in V}y_{i1}^2\le1,\ n^{-1}\sum_{i\in V}y_{i0}^2\le1\}\).
\par\smallskip\noindent\emph{Member properties.} \texttt{ass:\allowbreak{}ani}; \texttt{ass:\allowbreak{}treated-\allowbreak{}endpoint-\allowbreak{}l2}; \texttt{ass:\allowbreak{}control-\allowbreak{}endpoint-\allowbreak{}l2}.
\end{definition}

\begin{definition}[def:gate-conflict-graph]\label{def:gate-conflict-graph}
\par\smallskip\noindent\emph{Defined object.} \texttt{Two-\allowbreak{}copy GATE conflict graph}
\par\smallskip\noindent\emph{Construction.} \(H_G\) has vertex set \(\mathcal U\) and joins \((i,a)\) to \((j,b)\) if and only if \(a\ne b\) and \(\operatorname{dist}_G(i,j)\le2\).
\par\smallskip\noindent\emph{Inputs.} \texttt{G}.
\end{definition}

\begin{definition}[def:reservoir-capsule]\label{def:reservoir-capsule}
\par\smallskip\noindent\emph{Defined object.} \texttt{Balanced reservoir stable-\allowbreak{}set capsule}
\par\smallskip\noindent\emph{Construction.} Draw up to \(n^3\) uniform \(k\)-subsets \(A\); on the first draw with \(|F_G(A)|\ge k\), draw \(C\) uniformly from the \(k\)-subsets of \(F_G(A)\), swap the roles of \(A\) and \(C\) using \(\omega\), force the closed neighborhoods of \(T\) to treatment and those of \(C_0\) to control, and use fair independent assignments elsewhere; use one fair global arm if all trials fail.
\par\smallskip\noindent\emph{Inputs.} \texttt{G}; \texttt{k}.
\end{definition}

\begin{definition}[def:normalized-kernel]\label{def:normalized-kernel}
\par\smallskip\noindent\emph{Defined object.} \texttt{Signed normalized capsule covariance}
\par\smallskip\noindent\emph{Construction.} When every \(\pi_u>0\), \((\Gamma_G)_{(i,a),(j,b)}=s_as_b\{\Pr(X_{(i,a)}=X_{(j,b)}=1\mid G)/(\pi_{(i,a)}\pi_{(j,b)})-1\}\); the normalized kernel is undefined if any inclusion probability is zero.
\par\smallskip\noindent\emph{Inputs.} \texttt{\textbackslash{}\allowbreak{}mathsf D\_\allowbreak{}G\textasciicircum{}\{\textbackslash{}\allowbreak{}mathrm\{cap\}\allowbreak{}\}\allowbreak{}}.
\end{definition}

\begin{definition}[def:raw-estimator]\label{def:raw-estimator}
\par\smallskip\noindent\emph{Defined object.} \texttt{Focal Horvitz-\allowbreak{}-\allowbreak{}Thompson construction}
\par\smallskip\noindent\emph{Construction.} When every \(\pi_u>0\), \(\widetilde\tau_G=n^{-1}\sum_{i\in V}\sum_{a\in\{0,1\}}s_aX_{(i,a)}y_{ia}/\pi_{(i,a)}\); otherwise set \(\widetilde\tau_G=0\), consistently with the symbol-table convention.
\par\smallskip\noindent\emph{Inputs.} \texttt{\textbackslash{}\allowbreak{}mathsf D\_\allowbreak{}G\textasciicircum{}\{\textbackslash{}\allowbreak{}mathrm\{cap\}\allowbreak{}\}\allowbreak{}}; \texttt{y}.
\end{definition}

\begin{definition}[def:observed-estimator]\label{def:observed-estimator}
\par\smallskip\noindent\emph{Defined object.} \texttt{Observed-\allowbreak{}data capsule estimator}
\par\smallskip\noindent\emph{Construction.} \(\widehat\tau_G=\mathbb E[\widetilde\tau_G\mid G,Z^{\mathrm{cap}},y(Z^{\mathrm{cap}})]\).
\par\smallskip\noindent\emph{Inputs.} \texttt{\textbackslash{}\allowbreak{}widetilde\textbackslash{}\allowbreak{}tau\_\allowbreak{}G}; \texttt{Z\textasciicircum{}\{\textbackslash{}\allowbreak{}mathrm\{cap\}\allowbreak{}\}\allowbreak{}}.
\end{definition}

\begin{definition}[def:minimax-risk]\label{def:minimax-risk}
\par\smallskip\noindent\emph{Defined object.} Unrestricted endpoint-\(L^2\) full-GATE risk
\par\smallskip\noindent\emph{Construction.} \(R_2(G)=\inf_{(\mathsf D,T_*)}\sup_{y\in\mathcal M_2(G)}\mathbb E_{\mathsf D}[(T_*(G,Z,y(Z))-\tau_G(y))^2\mid G]\), with the infimum over every outcome-independent assignment law and every measurable possibly randomized estimator.
\par\smallskip\noindent\emph{Inputs.} \texttt{def:\allowbreak{}endpoint-\allowbreak{}class}.
\end{definition}

\begin{definition}[def:frontier-scale]\label{def:frontier-scale}
\par\smallskip\noindent\emph{Defined object.} \texttt{Random-\allowbreak{}regular stable-\allowbreak{}set frontier scale}
\par\smallskip\noindent\emph{Construction.} \(r(n,d)=(d+1)^2/[n\log(d+2)]\).
\end{definition}

\begin{definition}[def:local-compiler-handle]\label{def:local-compiler-handle}
\par\smallskip\noindent\emph{Defined object.} \texttt{Local covariance compiler handle}
\par\smallskip\noindent\emph{Construction.} Apply truncated radius-two avoidance messages to approximate the reservoir acceptance and inclusion probabilities, and use randomized matrix--vector probing under \(\mathsf D_G^{\mathrm{cap}}\) to target the action and norm of \(\Gamma_G\) together with the conditional Rao--Blackwell weights.
\par\smallskip\noindent\emph{Inputs.} \texttt{G}; \texttt{\textbackslash{}\allowbreak{}mathsf D\_\allowbreak{}G\textasciicircum{}\{\textbackslash{}\allowbreak{}mathrm\{cap\}\allowbreak{}\}\allowbreak{}}.
\end{definition}

\section{Main results}

\begin{proposition}[prop:capsule-realization]\label{prop:capsule-realization}
On every accepted capsule draw, \((T\times\{1\})\cup(C_0\times\{0\})\) is a stable set of \(H_G\), the closed neighborhoods forced by \(T\) and \(C_0\) are disjoint, and \(Z^{\mathrm{cap}}\) reveals \(y_{i1}\) for every \(i\in T\) and \(y_{i0}\) for every \(i\in C_0\).
\end{proposition}
\begin{proof}
Fix an accepted draw.  By the definition of \(F_G(A)\), every \(i\in A\) and \(j\in C\subseteq F_G(A)\) satisfy
\[
\operatorname{dist}_G(i,j)>2. \tag{1}
\]
The arm swap either gives \((T,C_0)=(A,C)\) or \((T,C_0)=(C,A)\), so (1) holds for every \(i\in T\) and \(j\in C_0\).  Edges of \(H_G\) join only opposite-arm vertices, and an opposite-arm pair \((i,1),(j,0)\) is an edge exactly when \(\operatorname{dist}_G(i,j)\le2\), by \(\text{def:gate-conflict-graph}\).  Hence no two vertices of
\[
(T\times\{1\})\cup(C_0\times\{0\})
\]
are adjacent, proving stability.

If a vertex \(v\) belonged to both \(\bigcup_{i\in T}N_G[i]\) and \(\bigcup_{j\in C_0}N_G[j]\), then for some \(i\in T\) and \(j\in C_0\),
\[
\operatorname{dist}_G(i,j)
\le \operatorname{dist}_G(i,v)+\operatorname{dist}_G(v,j)
\le2,
\]
contradicting (1).  Thus the two forced-neighborhood unions are disjoint, so the assignment in \(\text{def:reservoir-capsule}\) is well defined.  For \(i\in T\), it has \(Z^{\mathrm{cap}}|_{N_G[i]}=\boldsymbol 1|_{N_G[i]}\).  Assumption \(\text{ass:ani}\) therefore gives
\[
y_i(Z^{\mathrm{cap}})=y_i(\boldsymbol 1)=y_{i1}.
\]
The identical argument with zeros gives \(y_i(Z^{\mathrm{cap}})=y_{i0}\) for every \(i\in C_0\).  This proves all three assertions.
\end{proof}
\par\smallskip\noindent \textit{Justification.} This establishes that the internal focal sample corresponds to legal, simultaneously observed endpoint outcomes. \textit{Gap.} The maximal-stable-set realization in Kandiros, Harshaw, and Savje is general; this proposition specializes it to a balanced nonmaximal capsule and its explicit assignment map. \textit{Consumer.} It lets the estimator use two balanced endpoint samples without placing any restriction on irrelevant exposures.

\begin{lemma}[lem:exact-risk-certificate]\label{lem:exact-risk-certificate}
If \(\min_{u\in\mathcal U}\pi_u>0\), then \(\widetilde\tau_G\) and \(\widehat\tau_G\) are unbiased for \(\tau_G\), and for every \(y\in\mathcal M_2(G)\), \[\mathbb E[(\widehat\tau_G-\tau_G)^2\mid G]\le\mathbb E[(\widetilde\tau_G-\tau_G)^2\mid G]\le\frac{2\|\Gamma_G\|_{\mathrm{op}}}{n}.\]
\end{lemma}
\begin{proof}
Condition throughout on the revealed graph \(G\), and suppose that \(\pi_u>0\) for every \(u\in\mathcal U\).  All quantities below are integrable because \(\mathcal U\) is finite, the endpoint array is fixed and finite, and every denominator is strictly positive.  For \(u=(i,a)\), put \(W_u=X_u/\pi_u\).  The definition of \(\pi_u\) gives
\[
 \mathbb E[W_u\mid G]
 =\frac{\mathbb E[X_u\mid G]}{\pi_u}=1. \tag{1}
\]
On the event \(X_{(i,a)}=1\), Proposition \(\text{prop:capsule-realization}\) says that the realized assignment reveals the endpoint \(y_{ia}\).  Thus every summand entering the raw estimator is available on the event on which it is nonzero, and the conditional expectation in \(\text{def:observed-estimator}\) is an observed-data statistic.

By \(\text{def:raw-estimator}\), equation (1), linearity of expectation for the finite sum, the definition of \(s_a\), and the definition of \(\tau_G\),
\[
\begin{aligned}
 \mathbb E[\widetilde\tau_G\mid G]
 &=\frac1n\sum_{i\in V}\sum_{a\in\{0,1\}}
       s_a y_{ia}\mathbb E[W_{(i,a)}\mid G]\\
 &=\frac1n\sum_{i\in V}(y_{i1}-y_{i0})
 =\tau_G(y). \tag{2}
\end{aligned}
\]
Let \(Y\in\mathbb R^{\mathcal U}\) have coordinates \(Y_{(i,a)}=y_{ia}\), and let \(S\) be the diagonal matrix with \(S_{(i,a),(i,a)}=s_a\).  From \(\text{def:normalized-kernel}\), for \(u=(i,a)\) and \(v=(j,b)\),
\[
 (\Gamma_G)_{uv}
 =s_as_b\left\{
 \frac{\mathbb E[X_uX_v\mid G]}{\pi_u\pi_v}-1
 \right\}
 =s_as_b\operatorname{Cov}(W_u,W_v\mid G). \tag{3}
\]
Consequently
\[
 \Gamma_G=S\operatorname{Cov}\bigl((W_u)_{u\in\mathcal U}\mid G\bigr)S. \tag{4}
\]
The conditional covariance matrix in (4) is symmetric positive semidefinite, and congruence by \(S\) preserves positive semidefiniteness.  Hence \(\Gamma_G\) is symmetric positive semidefinite.  Expanding the variance of the finite sum in \(\text{def:raw-estimator}\) and applying (3) yields the exact identity
\[
 \operatorname{Var}(\widetilde\tau_G\mid G)
 =\frac1{n^2}\sum_{u,v\in\mathcal U}Y_u(\Gamma_G)_{uv}Y_v
 =\frac1{n^2}Y^{\mathsf T}\Gamma_GY. \tag{5}
\]
For \(y\in\mathcal M_2(G)\), the two membership conditions in \(\text{def:endpoint-class}\), equivalently \(\text{ass:treated-endpoint-l2}\) and \(\text{ass:control-endpoint-l2}\), imply
\[
 \lVert Y\rVert_2^2
 =\sum_{i\in V}y_{i1}^2+\sum_{i\in V}y_{i0}^2
 \le 2n. \tag{6}
\]
The Rayleigh-quotient inequality for the real symmetric matrix \(\Gamma_G\), followed by (6), therefore gives
\[
 \operatorname{Var}(\widetilde\tau_G\mid G)
 \le \frac{\lVert\Gamma_G\rVert_{\mathrm{op}}\lVert Y\rVert_2^2}{n^2}
 \le \frac{2\lVert\Gamma_G\rVert_{\mathrm{op}}}{n}. \tag{7}
\]

Let
\[
 \mathscr O=\sigma\bigl(G,Z^{\mathrm{cap}},y(Z^{\mathrm{cap}})\bigr).
\]
Because \(\sigma(G)\subseteq\mathscr O\), \(\text{def:observed-estimator}\), the tower property, and (2) imply
\[
 \mathbb E[\widehat\tau_G\mid G]
 =\mathbb E\!\left[\mathbb E[\widetilde\tau_G\mid\mathscr O]\mid G\right]
 =\mathbb E[\widetilde\tau_G\mid G]
 =\tau_G(y). \tag{8}
\]
Thus both estimators are conditionally unbiased.  The conditional variance decomposition, again using \(\widehat\tau_G=\mathbb E[\widetilde\tau_G\mid\mathscr O]\), gives
\[
\begin{aligned}
 \operatorname{Var}(\widetilde\tau_G\mid G)
 &=\operatorname{Var}(\widehat\tau_G\mid G)
   +\mathbb E\!\left[
      \operatorname{Var}(\widetilde\tau_G\mid\mathscr O)
      \,middle|\,G\right]\\
 &\ge \operatorname{Var}(\widehat\tau_G\mid G). \tag{9}
\end{aligned}
\]
Conditional mean squared error equals conditional variance for each estimator by (2) and (8).  Combining (7) and (9) proves
\[
 \mathbb E[(\widehat\tau_G-\tau_G)^2\mid G]
 \le \mathbb E[(\widetilde\tau_G-\tau_G)^2\mid G]
 \le \frac{2\lVert\Gamma_G\rVert_{\mathrm{op}}}{n},
\]
as claimed.
\end{proof}
\par\smallskip\noindent \textit{Justification.} The lemma turns the design problem into the single graph-dependent covariance-operator certificate used by the kernel theorem. \textit{Gap.} Theorem 4 and Corollary 5 of Harshaw, S\"avje, and Wang, arXiv:2210.08698v3 of 14 August 2025, provide the general Riesz-operator comparison, while the signed two-arm kernel and hidden-focal-set Rao--Blackwell step are specific to this design. \textit{Consumer.} Any improved bound on \(\Gamma_G\) immediately becomes a legal observed-data full-GATE risk guarantee.

\begin{openendedquestion}[oeq:capsule-stable-kernel]\label{oeq:capsule-stable-kernel}
Whenever all focal inclusion probabilities are positive, the capsule kernel has the exact decomposition
\[
\Gamma_G=
\begin{pmatrix}K_{11}&-K_{10}\\-K_{01}&K_{00}\end{pmatrix}\succeq0,
\qquad
(\Gamma_G)_{uu}=\pi_u^{-1}-1,
\qquad
\|\Gamma_G\|_{\mathrm{op}}\ge\max_u(\pi_u^{-1}-1),
\]
where \(K_{ab}(i,j)=\operatorname{Cov}(X_{(i,a)}/\pi_{(i,a)},X_{(j,b)}/\pi_{(j,b)}\mid G)\).  Its incompatible opposite-arm entries equal \(1\), and equality of the two focal arm totals makes the all-ones vector a null vector when the inclusion probabilities are uniform.  It remains a conjecture whether there are universal constants \(c,C>0\) such that, uniformly along every sequence with \(d\to\infty\) and \(d=o(\sqrt n)\), with probability \(1-o(1)\) over \(G_n\),
\[
p_{\mathrm{acc}}(G)\ge1-n^{-4},
\qquad
c\frac{k}{n}\le\min_{u\in\mathcal U}\pi_u
\le\max_{u\in\mathcal U}\pi_u\le C\frac{k}{n},
\qquad
\|\Gamma_G\|_{\mathrm{op}}\le C\frac{n}{k}.
\]
No \(O(n/k)\) capsule norm bound is asserted as proved.
\end{openendedquestion}
\par\smallskip\noindent \textit{Justification.} This conjecture preserves the exact covariance identities and isolates the one operator estimate that would be needed for a sharper capsule theorem. \textit{Gap.} No switching, mixing, or spectral argument currently controls the acceptance-conditioned compatible-pair block at order \(n/k\). \textit{Consumer.} It is not a premise of any proved headline; it directs future work on whether the capsule can improve the published upper rate.

\begin{proposition}[prop:finite-evaluator]\label{prop:finite-evaluator}
For every finite revealed graph \(G\), an exact draw from \(\mathsf D_G^{\mathrm{cap}}\) can be generated in randomized time polynomial in \(n\) and \(d\). Finite enumeration of the capped latent sample space computes every \(\pi_u\) and detects whether any \(\pi_u\) is zero. If every \(\pi_u>0\), the same enumeration computes every entry of \(\Gamma_G\), evaluates \(\widetilde\tau_G\) on every latent atom, and computes every conditional weight and value defining \(\widehat\tau_G\) on each positive-probability observed-data atom. Under this positivity condition, \(\widehat\tau_G\) is a finite measurable observed-data function and is exactly auditable without observing internal focal metadata. Exact enumeration is not asserted to run in polynomial time.
\end{proposition}
\begin{proof}
Fix a finite revealed simple \(d\)-regular graph \(G\). A single proposal draws a uniform \(k\)-subset \(A\), computes \(F_G(A)\) by exploring graph distance at most two from each vertex of \(A\), and tests whether \(\lvert F_G(A)\rvert\ge k\). Each depth-two exploration examines at most
\[
1+d+d(d-1)
\]
vertex incidences. Since \(k\le n\), one proposal uses a number of elementary operations polynomial in \(n\) and \(d\). There are at most \(n^3\) proposals. If one is accepted, drawing a uniform \(k\)-subset \(C\subseteq F_G(A)\), drawing \(\omega\), forming the two forced unions of closed neighborhoods, and drawing at most \(n\) remaining fair assignment coins are all finite operations taking time polynomial in \(n\) and \(d\). If no proposal is accepted, only the fair global coin is drawn. Thus the construction in \(\text{def:reservoir-capsule}\) gives an exact draw from \(\mathsf D_G^{\mathrm{cap}}\) in randomized time polynomial in \(n\) and \(d\).

For the exact finite audit, let \(\Omega_G\) be the latent sample space whose atoms record the ordered reservoir proposals up to the first acceptance or the cap of \(n^3\), the subsequent \(C\) when acceptance occurs, \(\omega\), and all assignment coins used by the construction. Every recorded choice ranges over a finite set, so \(\Omega_G\) is finite. For each \(\eta\in\Omega_G\), the construction specifies an explicit rational probability mass \(p_G(\eta)\), and the maps
\[
\eta\longmapsto Z^{\mathrm{cap}}(\eta)
\qquad\text{and}\qquad
\eta\longmapsto X_u(\eta)
\]
are deterministic. Hence, for every \(u\in\mathcal U\), finite summation gives
\[
\pi_u
=\mathbb E_{\mathsf D_G^{\mathrm{cap}}}[X_u\mid G]
=\sum_{\eta\in\Omega_G}p_G(\eta)X_u(\eta).
\tag{1}
\]
Every summand in (1) is a nonnegative rational number. Exact rational arithmetic therefore computes \(\pi_u\) and decides whether \(\pi_u=0\). Repeating (1) for the \(2n\) vertices \(u\in\mathcal U\) detects whether any focal inclusion probability vanishes.

Assume for the remainder of the evaluation argument that \(\pi_u>0\) for every \(u\in\mathcal U\). For \(u,v\in\mathcal U\), the same enumeration computes
\[
q_{uv}
:=\Pr(X_u=X_v=1\mid G)
=\sum_{\eta\in\Omega_G}p_G(\eta)X_u(\eta)X_v(\eta).
\tag{2}
\]
The denominators \(\pi_u\pi_v\) are nonzero by the displayed positivity condition. Substitution of (1) and (2) into \(\text{def:normalized-kernel}\) therefore computes every entry of \(\Gamma_G\) exactly.

On every accepted atom, \(\text{prop:capsule-realization}\) states that each endpoint outcome multiplied by a nonzero focal indicator in \(\widetilde\tau_G\) is revealed by \((Z^{\mathrm{cap}},y(Z^{\mathrm{cap}}))\). On a fallback atom, \(X_u=0\) for every \(u\) by the definition of the focal indicators. Consequently \(\text{def:raw-estimator}\), together with the positive probabilities computed in (1), evaluates \(\widetilde\tau_G(\eta)\) exactly on every latent atom: an accepted atom uses only its focal metadata and endpoint outcomes revealed under its assignment, and a fallback atom has value zero.

Fix an assignment \(z\in\{0,1\}^V\) having positive probability, and define its finite latent fiber by
\[
\Omega_G(z)
:=\bigl\{\eta\in\Omega_G:Z^{\mathrm{cap}}(\eta)=z\bigr\}.
\]
Because the potential-outcome schedule is fixed, \(y(Z^{\mathrm{cap}})\) is a deterministic function of \(Z^{\mathrm{cap}}\). Conditioning on
\[
(G,Z^{\mathrm{cap}},y(Z^{\mathrm{cap}}))=(G,z,y(z))
\]
therefore selects the same latent fiber as conditioning on \((G,Z^{\mathrm{cap}})=(G,z)\). Its probability
\[
P_G(z):=\sum_{\eta\in\Omega_G(z)}p_G(\eta)
\]
is strictly positive by the choice of \(z\). The conditional expectation in \(\text{def:observed-estimator}\) is thus the explicit finite ratio
\[
\widehat\tau_G(G,z,y(z))
=\frac{\displaystyle\sum_{\eta\in\Omega_G(z)}p_G(\eta)\widetilde\tau_G(\eta)}
{\displaystyle P_G(z)}.
\tag{3}
\]
Equation (3) computes both the conditional weights \(p_G(\eta)/P_G(z)\) and the value of \(\widehat\tau_G\). By \(\text{prop:capsule-realization}\), every outcome required by every summand in the numerator is a coordinate of the observed vector \(y(z)\); the internal focal metadata are summed out over the fiber. The positive-probability observed-data support is finite, so the function specified by (3) is measurable. This proves the claimed exact auditability from observed data without observing the internal focal metadata. Nothing in this finite-sum argument bounds \(\lvert\Omega_G\rvert\) by a polynomial, so it makes no polynomial-time claim for exact enumeration.

Finally, positivity is not automatic, even at the smallest admissible population size. Let \(n=4\), \(d=3\), and \(G=K_4\). Then \(k=1\), while every vertex is within graph distance two of every singleton \(A\), so \(F_G(A)=\varnothing\). All \(n^3\) proposals fail, the fallback is used, and \(X_u=0\) for every \(u\in\mathcal U\). Equation (1) then gives \(\pi_u=0\) for every \(u\), which the enumeration detects. The divisions in the positive-inclusion formulas for \(\Gamma_G\) and the focal Horvitz--Thompson expression are consequently unavailable on this example. This verifies that the positivity qualification used above is substantive and completes the proof.
\end{proof}
\par\smallskip\noindent \textit{Justification.} The finite experiment must expose zero inclusion before any division-based kernel or estimator formula is used. \textit{Gap.} Enumeration is exact but may be exponential, and positive focal inclusion is not automatic on every finite graph. \textit{Consumer.} It gives an exact small-instance audit and a benchmark for scalable approximations, conditional on verified positive inclusion when the kernel or estimators are evaluated.

\begin{theorem}[thm:published-conflict-graph-upper]\label{thm:published-conflict-graph-upper}
There is a universal constant \(C<\infty\) such that, uniformly along every sequence satisfying \(d\to\infty\) and \(d=o(\sqrt n)\), for every sufficiently large \(n\) and every revealed simple \(d\)-regular graph \(G\), the published Conflict Graph Design followed by observed-data Rao--Blackwellization gives an outcome-independent design and measurable estimator \(T_{\mathrm{CG}}^{\mathrm{RB}}\) satisfying \[\sup_{y\in\mathcal M_2(G)}\mathbb E[(T_{\mathrm{CG}}^{\mathrm{RB}}-\tau_G)^2\mid G]\le 25\lambda(H_G)/n\le C d^2/n.\] Consequently \(R_2(G)\le C d^2/n\). This theorem does not assert \(O(d^2/[n\log(d+2)])\) for the capsule.
\end{theorem}
\begin{proof}
Fix a sufficiently large index \(n\), a revealed simple \(d\)-regular graph \(G\), and a fixed schedule \(y\in\mathcal M_2(G)\).  Let \(T_{\mathrm{CG}}\) denote the modified Horvitz--Thompson estimator computed with the internal desired-exposure variables of the published Conflict Graph Design.  The cited result \(\text{lem:published-conflict-graph-risk}\) gives
\[
\mathbb E(T_{\mathrm{CG}}\mid G)=\tau_G(y) \tag{9}
\]
and
\[
\operatorname{Var}(T_{\mathrm{CG}}\mid G)
\le \frac{12.5\|B_G\|_{\mathrm{op}}}{n}
\left\{\frac1n\sum_{i\in V}y_{i1}^2+
\frac1n\sum_{i\in V}y_{i0}^2\right\}. \tag{10}
\]
The frozen minimax class permits estimators based on \((G,Z,y(Z))\), not hidden design metadata.  Define the observed-data version
\[
T_{\mathrm{CG}}^{\mathrm{RB}}
=\mathbb E[T_{\mathrm{CG}}\mid G,Z,y(Z)]. \tag{11}
\]
This is one common, data-driven map rather than a schedule-specific oracle.  Indeed, all spaces are finite, the joint law of the internal desired-exposure variables \(U\) and \(Z\) is fixed by \(G\), and \(T_{\mathrm{CG}}\) depends on the schedule only through \(y(Z)\).  Thus, for every realized pair \((z,Y)\), (11) is the same design-known finite sum
\[
T_{\mathrm{CG}}^{\mathrm{RB}}(G,z,Y)
=\sum_u T_{\mathrm{CG}}(G,z,Y,u)\Pr_G\{U=u\mid Z=z\},
\]
with an arbitrary value assigned on zero-probability assignments.  Hence it is a measurable function of the permitted observed data for every schedule in the class.  The tower property and (9) imply
\[
\mathbb E(T_{\mathrm{CG}}^{\mathrm{RB}}\mid G)
=\mathbb E(T_{\mathrm{CG}}\mid G)=\tau_G(y). \tag{12}
\]
The law of total variance gives
\[
\operatorname{Var}(T_{\mathrm{CG}}\mid G)
=\operatorname{Var}(T_{\mathrm{CG}}^{\mathrm{RB}}\mid G)
+\mathbb E\!\left[\operatorname{Var}(T_{\mathrm{CG}}\mid G,Z,y(Z))\mid G\right], \tag{13}
\]
so the second nonnegative term shows that Rao--Blackwellization cannot increase variance.

By \(\text{def:endpoint-class}\), the two endpoint second moments in braces in (10) are separately at most one.  Combining (10)--(13) and using unbiasedness yields, uniformly over \(y\in\mathcal M_2(G)\),
\[
\mathbb E[(T_{\mathrm{CG}}^{\mathrm{RB}}-\tau_G)^2\mid G]
=\operatorname{Var}(T_{\mathrm{CG}}^{\mathrm{RB}}\mid G)
\le\frac{25\|B_G\|_{\mathrm{op}}}{n}. \tag{14}
\]
The proved spectral identification \(\text{lem:gate-conflict-spectral-map}\) changes (14) into
\[
\sup_{y\in\mathcal M_2(G)}
\mathbb E[(T_{\mathrm{CG}}^{\mathrm{RB}}-\tau_G)^2\mid G]
\le\frac{25\lambda(H_G)}{n}
\le\frac{25(1+d^2)}{n}. \tag{15}
\]
Assumption \(\text{ass:degree-diverges}\) implies that eventually \(d\ge1\), and then \(1+d^2\le2d^2\).  Thus (15) is at most \(50d^2/n\).

The Conflict Graph Design depends only on \(G\), and (11) is an observed-data measurable estimator; hence this pair is admissible in the infimum of \(\text{def:minimax-risk}\).  Evaluating that infimum at this one pair and applying (15) proves
\[
R_2(G)\le
\sup_{y\in\mathcal M_2(G)}
\mathbb E[(T_{\mathrm{CG}}^{\mathrm{RB}}-\tau_G)^2\mid G]
\le\frac{50d^2}{n}. \tag{16}
\]
The bound is deterministic for every revealed simple \(d\)-regular \(G\) once \(d\ge1\); therefore it holds, a fortiori, with probability \(1-o(1)\) under \(G_n\).  Assumption \(\text{ass:subsqrt-degree}\) specifies the common regime in which this upper bound is paired with the frozen lower theorem, although it is not needed for the finite-graph inequality (16).
\end{proof}
\par\smallskip\noindent \textit{Justification.} The retargeted theorem proves the valid upper side by Rao--Blackwellizing the published Conflict Graph Design estimator into the frozen observed-data minimax class. \textit{Gap.} Its order is \(d^2/n\), leaving a logarithmic factor above the all-procedure lower edge. \textit{Consumer.} The logarithmic-gap bracket theorem uses it as the upper side.

\begin{theorem}[thm:all-procedure-lower]\label{thm:all-procedure-lower}
There is a universal constant \(c>0\) such that, for every feasible degree sequence satisfying \(d=o(\sqrt n)\), \[\Pr_{G_n}\!\left\{R_2(G)\ge c\frac{1+d}{n}\right\}\longrightarrow1.\] For every revealed simple graph \(G\), one also has the deterministic branch \(R_2(G)\ge1/(16n)\). These bounds hold simultaneously against every outcome-independent assignment law and every measurable, possibly biased or randomized estimator on \(\mathcal M_2(G)\). The stronger lower edge of order \(r(n,d)\) is not asserted.
\end{theorem}
\begin{proof}
The fixed-graph certificate \(\text{lem:balanced-hole-bayes-lower}\) gives, for every revealed graph and without an asymptotic condition,
\[
R_2(G)\ge\frac1{16n}. \tag{16}
\]
This conclusion already ranges over all outcome-independent designs and all measurable, possibly biased or randomized estimators because those are the procedures in \(\text{def:minimax-risk}\), and the proof of the certificate took the infimum over that entire class.

For the growing-degree contribution, Corollary 5.4 of Kandiros, Harshaw, and S\"avje, materialized as \(\text{lem:published-random-regular-gate-lower}\), states that there is a universal \(c_0>0\) and events \(L_n\), measurable with respect to \(G_n\), such that
\[
\Pr_{G_n}(L_n)\longrightarrow1,
\qquad
R_2(G)\ge c_0\frac d n\quad\text{on }L_n. \tag{17}
\]
The notation agrees exactly: the cited paper's arbitrary-neighborhood-interference class at moment parameter \(q=2\) imposes, separately on the two GATE endpoint arms, the two empirical second-moment inequalities in \(\text{def:endpoint-class}\); its GATE is the all-treated minus all-control contrast \(\tau_G\); and its minimax risk takes the infimum over outcome-independent designs and estimators, which is \(R_2(G)\) in \(\text{def:minimax-risk}\).  If an estimator additionally uses an independent random seed \(U\), then
\[
\overline T(G,Z,y(Z))
=\mathbb E[T(G,Z,y(Z),U)\mid G,Z,y(Z)]
\]
is observed-data measurable and, for each fixed schedule \(y\), conditional Jensen gives
\[
\mathbb E[(\overline T-\tau_G(y))^2\mid G]
\le
\mathbb E[(T-\tau_G(y))^2\mid G].
\]
Thus allowing estimator randomization cannot reduce the minimax risk, and no superpopulation randomness or estimator restriction is introduced in passing from the cited corollary to (17).

On \(L_n\), equations (16)--(17) give
\[
R_2(G)\ge
\max\!\left\{\frac1{16n},c_0\frac d n\right\}
\ge \frac12\left\{\frac1{16n}+c_0\frac d n\right\}
\ge c\frac{1+d}{n}, \tag{18}
\]
where \(c=\tfrac12\min\{1/16,c_0\}>0\) is universal.  Taking probabilities in (18) and using \(\Pr_{G_n}(L_n)\to1\) proves the displayed assertion for every feasible sequence satisfying \(\text{ass:subsqrt-degree}\).

The earlier proposed improvement would require the additional random-graph conclusion
\[
b(G)=O_{\Pr}\!\left(\frac{n\log(d+2)}{(d+1)^2}\right). \tag{19}
\]
Indeed, substituting (19) into (15) would produce the scale \(r(n,d)\).  The previously supplied switching sketch did not define the exposure-conditioned residual degree sequences or switching strata and did not establish uniform valid-to-invalid forward and reverse counts after excluding loops, repeated edges, exceptional internal-edge levels, and shortages of available half-edges.  Corollary 5.4 supplies (17), not (19).  Therefore (19) is not used here, and the theorem honestly stops at (18).
\end{proof}
\par\smallskip\noindent \textit{Justification.} A stable-set upper bound is not an optimal-design theorem without an all-procedure converse on the same endpoint class. \textit{Gap.} The proved growing-degree lower edge is the published order \(d/n\). Raising it to \(r(n,d)\) requires the still-open uniform random-regular balanced-hole estimate. \textit{Consumer.} The converse rules out risk \(o((1+d)/n)\) for every cluster, conflict-priority, nonlinear, biased, or randomized procedure in the frozen class.

\begin{proposition}[prop:no-interference-reduction]\label{prop:no-interference-reduction}
When \(d=0\), the capsule accepts immediately, selects two disjoint uniform fixed-size samples, labels them by the arm-swap coin, and assigns independent fair coins to the remaining units.  Its raw estimator is the difference of the corresponding endpoint sample means, its Rao--Blackwellization has worst-case risk \(O(n^{-1})\), and the unrestricted minimax risk satisfies \(R_2(G)=\Theta(n^{-1})\).
\end{proposition}
\begin{proof}
Fix an index \(n\) and suppose \(d=0\).  A simple \(0\)-regular graph has no edges, so \(N_G[i]=\{i\}\), distinct vertices have infinite graph distance, and therefore
\[
 F_G(A)=V\setminus A. \tag{1}
\]
Because the definition of \(k\) ensures \(k\le\lfloor n/4\rfloor\), equation (1) gives \(|F_G(A)|=n-k\ge k\).  Hence the first trial is accepted under \(\text{def:reservoir-capsule}\).

Conditional on \(A\), the set \(C\) is uniform among the \(k\)-subsets of \(V\setminus A\).  For any fixed \(k\)-subset \(S\subseteq V\),
\[
\begin{aligned}
 \Pr(C=S\mid G)
 &=\sum_{\substack{A\subseteq V:\ |A|=k,\ A\cap S=\varnothing}}
   \frac1{\binom nk}\frac1{\binom{n-k}{k}}\\
 &=\frac{\binom{n-k}{k}}
         {\binom nk\binom{n-k}{k}}
 =\frac1{\binom nk}. \tag{2}
\end{aligned}
\]
Thus \(A\) and \(C\) are marginally uniform \(k\)-subsets and are disjoint.  The independent fair arm-swap coin makes each of \(T\) and \(C_0\) marginally uniform among the \(k\)-subsets of \(V\), while \(T\cap C_0=\varnothing\).  Since closed neighborhoods are singletons, \(\text{def:reservoir-capsule}\) treats every member of \(T\), controls every member of \(C_0\), and assigns mutually independent fair coins to the remaining units.  This proves the asserted design reduction.

For every \(i\in V\) and \(a\in\{0,1\}\), marginal uniformity and the arm swap give
\[
 \pi_{(i,a)}=\frac{k}{n}>0. \tag{3}
\]
Proposition \(\text{prop:capsule-realization}\) supplies the corresponding observed endpoint whenever the focal indicator is one.  Substitution of (3) into \(\text{def:raw-estimator}\) therefore gives
\[
 \widetilde\tau_G
 =\frac1k\sum_{i\in T}y_{i1}
  -\frac1k\sum_{i\in C_0}y_{i0}. \tag{4}
\]
This is the difference of the two endpoint sample means stated in the proposition.

We next derive a uniform risk bound.  If \(S\) is a uniform \(k\)-subset of \(V\), write \(I_i=\mathbf 1\{i\in S\}\), \(\bar x=n^{-1}\sum_i x_i\), and \(\bar x_S=k^{-1}\sum_i I_ix_i\).  Direct counting gives
\[
 \mathbb E[I_i]=\frac{k}{n},\qquad
 \mathbb E[I_iI_j]=\frac{k(k-1)}{n(n-1)}\quad(i\ne j). \tag{5}
\]
Expanding the first two moments using (5) yields
\[
 \mathbb E[\bar x_S]=\bar x,
 \qquad
 \operatorname{Var}(\bar x_S)
 =\frac{n-k}{k(n-1)}
   \left\{\frac1n\sum_{i\in V}(x_i-\bar x)^2\right\}. \tag{6}
\]
Apply (6) separately to the marginal samples \(T\) and \(C_0\).  For every \(y\in\mathcal M_2(G)\), \(\text{def:endpoint-class}\), \(\text{ass:treated-endpoint-l2}\), and \(\text{ass:control-endpoint-l2}\) give, for \(a\in\{0,1\}\),
\[
 \frac1n\sum_{i\in V}(y_{ia}-\bar y_a)^2
 =\frac1n\sum_{i\in V}y_{ia}^2-\bar y_a^2
 \le 1. \tag{7}
\]
In particular, (4)--(6) show that \(\widetilde\tau_G\) is unbiased for \(\tau_G\).  If \(U=k^{-1}\sum_{i\in T}y_{i1}\) and \(W=k^{-1}\sum_{i\in C_0}y_{i0}\), then Cauchy--Schwarz implies
\[
 \operatorname{Var}(U-W)
 \le 2\operatorname{Var}(U)+2\operatorname{Var}(W).
\]
Together with (6)--(7), this proves the explicit bound
\[
 \sup_{y\in\mathcal M_2(G)}
 \mathbb E[(\widetilde\tau_G-\tau_G)^2\mid G]
 \le \frac{4(n-k)}{k(n-1)}. \tag{8}
\]
Condition (3) permits application of \(\text{lem:exact-risk-certificate}\), so \(\widehat\tau_G\) is unbiased and its mean squared error is no larger than the left side of (8).  From the definition of \(k\),
\[
 \frac{k}{n}\longrightarrow\frac{\log 2}{32}, \tag{9}
\]
because \(0<\log 2/32<1/4\), the lower truncation by one is eventually inactive, and division by \(n\) removes the floor error.  Hence (8)--(9) show that the observed-data Rao--Blackwell estimator has worst-case risk \(O(n^{-1})\).  The capsule law is outcome-independent, so it and \(\widehat\tau_G\) form an admissible pair in \(\text{def:minimax-risk}\).  Thus, for an absolute finite constant \(C\) and all sufficiently large \(n\),
\[
 R_2(G)\le \frac Cn. \tag{10}
\]

For the matching order lower bound, construct a prior on admissible fixed schedules as follows.  Let \(\{\varepsilon_{ia}:i\in V,\ a\in\{0,1\}\}\) be mutually independent Rademacher variables, independent of every outcome-independent design, and set
\[
 y_i(z)=\varepsilon_{i,z_i}. \tag{11}
\]
If a triangular-array realization is required by \(\text{def:endpoint-class}\), take such independent signs at every index.  Since \(N_G[i]=\{i\}\), equality of \(z\) and \(z'\) on \(N_G[i]\) implies \(z_i=z_i'\); hence (11) satisfies \(\text{ass:ani}\).  Moreover,
\[
 \frac1n\sum_{i\in V}y_{ia}^2=1,
 \qquad a\in\{0,1\}, \tag{12}
\]
so every realization belongs to \(\mathcal M_2(G)\) by \(\text{def:endpoint-class}\).

Fix any outcome-independent assignment law \(\mathsf D\).  Conditional on a realized assignment \(Z=z\) and the observed outcomes \((\varepsilon_{i,z_i})_{i\in V}\), independence of \(\mathsf D\) from the fixed schedule implies that the unobserved signs \((\varepsilon_{i,1-z_i})_{i\in V}\) remain mutually independent Rademacher variables.  The causal estimand has the decomposition
\[
 \tau_G
 =\frac1n\sum_{i\in V}
   \left\{(2z_i-1)\varepsilon_{i,z_i}
   +(1-2z_i)\varepsilon_{i,1-z_i}\right\}. \tag{13}
\]
The first term within braces is observed.  The second has conditional mean zero, its summands are conditionally independent, and \((1-2z_i)^2=1\).  Therefore, with \(\mathscr O=\sigma(G,Z,y(Z))\),
\[
 \operatorname{Var}(\tau_G\mid\mathscr O)
 =\frac1{n^2}\sum_{i\in V}(1-2Z_i)^2
 =\frac1n
 \quad\text{almost surely}. \tag{14}
\]
For any measurable, possibly randomized estimator \(T_*\), adjoin its randomization seed to \(\mathscr O\).  That seed is independent of the unobserved signs.  Conditional squared-error decomposition, or equivalently the optimality of conditional expectation under squared loss, then gives
\[
 \mathbb E[(T_* -\tau_G)^2]
 \ge \mathbb E[\operatorname{Var}(\tau_G\mid\mathscr O)]
 =\frac1n. \tag{15}
\]
The prior is supported on \(\mathcal M_2(G)\) by (11)--(12).  Hence the supremum risk over \(y\in\mathcal M_2(G)\) of the fixed pair \((\mathsf D,T_*)\) is at least its prior-average risk in (15).  Taking the infimum over all eligible pairs in \(\text{def:minimax-risk}\) yields
\[
 R_2(G)\ge\frac1n. \tag{16}
\]
Combining (10) and (16) proves \(R_2(G)=\Theta(n^{-1})\) when \(d=0\).

Finally, the smallest admissible instance supplies an exact audit.  If \(n=4\), then the definition gives \(k=1\).  There are twelve ordered pairs \((T,C_0)=(\{t\},\{c\})\) with \(t\ne c\), each of probability \(1/12\): each orientation receives probability \(1/24\) before summing the two arm-swap histories.  Conditional on each ordered pair, the other two assignments have four equiprobable fair-coin configurations.  Equation (3) gives \(\pi_{(i,a)}=1/4\) for all eight focal exposures, so there is no zero-inclusion case.  On every one of these design histories, (4) is \(y_{t1}-y_{c0}\), and exhaustive averaging gives
\[
 \frac1{12}\sum_{t\ne c}(y_{t1}-y_{c0})
 =\frac14\sum_{i=1}^4(y_{i1}-y_{i0})
 =\tau_G. \tag{17}
\]
Thus no enumerated assignment history contradicts unbiasedness.  Under the supported prior (11), equation (14) equals \(1/4\) for every one of the assignments, exactly matching the finite-instance lower-bound calculation.  The treated count is \(1+\operatorname{Binomial}(2,1/2)\), so it takes values \(1,2,3\); this also verifies that the capsule reduces to the stated two-stage randomization, not to fixed balanced complete randomization.
\end{proof}
\par\smallskip\noindent \textit{Justification.} This verifies that the construction reduces to an explicit no-interference two-stage design when radius-two conflicts disappear. \textit{Gap.} The reduction is a required sanity check rather than a novelty claim. \textit{Consumer.} It supplies the exact no-interference benchmark for interpreting the wider degree-dependent bracket without claiming a general exact frontier.

\begin{openendedquestion}[oeq:local-compiler]\label{oeq:local-compiler}
Can the local handle \(\mathsf H_{\mathrm{loc}}\) be developed into an algorithm running in time polynomial in \(n\) and \(d\) that evaluates or approximates \(\widehat\tau_G\) from \((G,Z^{\mathrm{cap}},y(Z^{\mathrm{cap}}))\), and returns an observable certificate \(U_G\) satisfying \[\sup_{y\in\mathcal M_2(G)}\mathbb E[(\widehat\tau_G-\tau_G)^2\mid G]\le U_G?\] Can such a certificate either establish \(U_G=O(r(n,d))\) on a probability-\(1-o(1)\) random-regular graph event, or furnish a rigorous obstruction to \(\|\Gamma_G\|_{\mathrm{op}}=O(n/k)\), while numerical approximation error contributes \(o(r(n,d))\)? If a balanced-hole lower edge at scale \(r(n,d)\) is separately established, an affirmative compiler and \(O(n/k)\) certificate could potentially close the resulting logarithmic bracket for \(R_2(G)\) and make the capsule deployable; it would not implement an already established matched minimax theorem. Both the polynomial-time algorithmic construction and the capsule spectral bound remain open, and the proved upper bound uses a different published Conflict Graph Design.
\end{openendedquestion}
\par\smallskip\noindent \textit{Justification.} Exact enumeration proves finite evaluability on the positive-inclusion domain but does not deliver a scalable estimator or risk certificate for large networks. \textit{Gap.} No inverse-polynomial overlap or observed-assignment-fiber bound, approximate-counting theorem, or operator-norm certification argument is available from the frozen primitives. \textit{Consumer.} If the balanced-hole lower edge is first established, a successful compiler together with an \(O(n/k)\) capsule certificate could potentially close the resulting logarithmic bracket and make the capsule deployable; it would not implement an already matched minimax theorem.

\begin{lemma}[lem:capsule-kernel-identities]\label{lem:capsule-kernel-identities}
Suppose that \(\pi_u>0\) for every \(u\in\mathcal U\).  For \(a,b\in\{0,1\}\), define the \(V\)-by-\(V\) block
\[
K_{ab}(i,j)=\operatorname{Cov}\!\left(
 \frac{X_{(i,a)}}{\pi_{(i,a)}},
 \frac{X_{(j,b)}}{\pi_{(j,b)}}\,\middle|\,G\right).
\]
Then \(\Gamma_G\) is positive semidefinite and, in the arm order \((1,0)\),
\[
\Gamma_G=
\begin{pmatrix}
K_{11}&-K_{10}\\
-K_{01}&K_{00}
\end{pmatrix},
\qquad
(\Gamma_G)_{uu}=\pi_u^{-1}-1,
\qquad
\|\Gamma_G\|_{\mathrm{op}}
 \ge \max_{u\in\mathcal U}(\pi_u^{-1}-1).
\]
If \(a\ne b\) and \(\operatorname{dist}_G(i,j)\le2\), then
\((\Gamma_G)_{(i,a),(j,b)}=1\).  Moreover, on every realization of the capped design,
\[
\sum_{i\in V}X_{(i,1)}=\sum_{i\in V}X_{(i,0)}.
\]
Consequently, if all \(\pi_u\) have one common positive value, the all-ones vector belongs to the kernel of \(\Gamma_G\).
\end{lemma}
\begin{proof}
Condition throughout on the revealed graph \(G\), and write all expectations and covariances with respect to \(\mathsf D_G^{\mathrm{cap}}\) conditional on \(G\).  The hypothesis \(\pi_u>0\) makes every following normalization well defined.  Define
\[
 W_u=\frac{X_u}{\pi_u},
 \qquad
 D_{uv}=\mathbf 1\{u=v\}s_a
 \quad\text{for }u=(i,a),
\]
so that \(D\) is diagonal and \(D^2=I\).  Since \(\pi_u=\mathbb E[X_u\mid G]\),
\[
 \mathbb E[W_u\mid G]=1. \tag{1}
\]
Consequently, the probability formula in \(\text{def:normalized-kernel}\) is equivalent, entry by entry, to
\[
 \Gamma_G=D\operatorname{Cov}(W\mid G)D. \tag{2}
\]
Indeed, the \((u,v)\) entry on the right of (2) is
\[
 s_as_b\left\{
 \frac{\mathbb E[X_uX_v\mid G]}{\pi_u\pi_v}-1
 \right\}, \tag{3}
\]
where (1) supplies the subtracted term.

For any deterministic \(q\in\mathbb R^{\mathcal U}\), covariance bilinearity and (2) give
\[
 q^{\mathsf T}\Gamma_Gq
 =(Dq)^{\mathsf T}\operatorname{Cov}(W\mid G)(Dq)
 =\operatorname{Var}\!\left((Dq)^{\mathsf T}W\mid G\right)
 \ge 0. \tag{4}
\]
Thus \(\Gamma_G\) is positive semidefinite.  Splitting (2) into its arm blocks and using \(s_1=1\) and \(s_0=-1\) yields, in arm order \((1,0)\),
\[
 \Gamma_G=
 \begin{pmatrix}
 K_{11}&-K_{10}\\
 -K_{01}&K_{00}
 \end{pmatrix}. \tag{5}
\]
For \(u=(i,a)\), the indicator \(X_u\) is Bernoulli with conditional mean \(\pi_u\).  Since \(s_a^2=1\), (2) therefore gives
\[
 (\Gamma_G)_{uu}
 =\operatorname{Var}\!\left(\frac{X_u}{\pi_u}\,\middle|\,G\right)
 =\frac{\pi_u(1-\pi_u)}{\pi_u^2}
 =\pi_u^{-1}-1. \tag{6}
\]
Let \(e_u\) be the Euclidean coordinate vector at \(u\).  By the definition of the operator norm and the Cauchy--Schwarz inequality,
\[
 \|\Gamma_G\|_{\mathrm{op}}
 \ge \|\Gamma_Ge_u\|_2
 \ge \left|e_u^{\mathsf T}\Gamma_Ge_u\right|
 =\pi_u^{-1}-1, \tag{7}
\]
where the last quantity is nonnegative by \(0<\pi_u\le1\).  Maximizing (7) over \(u\in\mathcal U\) proves the asserted operator-norm lower bound.

Now fix \(a\ne b\) and units \(i,j\) satisfying \(\operatorname{dist}_G(i,j)\le2\).  On an accepted trial, \(C\subseteq F_G(A)\) by \(\text{def:reservoir-capsule}\).  Hence every \(h\in A\) and \(\ell\in C\) satisfy
\[
 \operatorname{dist}_G(h,\ell)>2. \tag{8}
\]
The arm-swap coin only exchanges the roles of \(A\) and \(C\), so (8) implies that \(i\) and \(j\) cannot simultaneously belong to the two opposite-arm focal sets.  Thus
\[
 X_{(i,a)}X_{(j,b)}=0 \tag{9}
\]
on every accepted realization.  On the capped fallback all focal indicators equal zero by \(\text{def:reservoir-capsule}\), so (9) also holds there.  Positivity of the two marginal inclusion probabilities permits division; using (1) and (9),
\[
 \operatorname{Cov}\!\left(
 \frac{X_{(i,a)}}{\pi_{(i,a)}},
 \frac{X_{(j,b)}}{\pi_{(j,b)}}\,\middle|\,G
 \right)
 =\frac{0}{\pi_{(i,a)}\pi_{(j,b)}}-1
 =-1. \tag{10}
\]
Because \(a\ne b\) entails \(s_as_b=-1\), equations (2) and (10) imply
\[
 (\Gamma_G)_{(i,a),(j,b)}=1. \tag{11}
\]

On every accepted realization, the construction in \(\text{def:reservoir-capsule}\) has \(|A|=|C|=k\), and the swap gives \(|T|=|C_0|=k\).  On every fallback realization both focal sets are empty.  Therefore, pathwise,
\[
 \sum_{i\in V}X_{(i,1)}
 =\sum_{i\in V}X_{(i,0)}. \tag{12}
\]
Suppose finally that \(\pi_u=\pi>0\) for all \(u\in\mathcal U\), and let \(\mathbf 1\) denote the all-ones vector in \(\mathbb R^{\mathcal U}\).  Equation (12) gives the pathwise identity
\[
 (D\mathbf 1)^{\mathsf T}W
 =\frac{1}{\pi}
 \left\{
 \sum_{i\in V}X_{(i,1)}-\sum_{i\in V}X_{(i,0)}
 \right\}
 =0. \tag{13}
\]
Taking the covariance of every coordinate of \(W\) with the scalar in (13), then using (2), yields
\[
 \Gamma_G\mathbf 1
 =D\operatorname{Cov}(W\mid G)D\mathbf 1
 =D\operatorname{Cov}\!\left(W,(D\mathbf 1)^{\mathsf T}W\mid G\right)
 =0. \tag{14}
\]
Thus the all-ones vector belongs to the kernel of \(\Gamma_G\).

For an exact smallest-instance stress test, take \(n=4\) and \(d=0\).  Then \(k=1\), every first reservoir is accepted, and the final ordered pair \((T,C_0)\) is uniform over the twelve ordered pairs of distinct units.  Hence \(\pi_u=1/4\) for every \(u\).  Direct enumeration gives diagonal entries \(3\), distinct same-arm entries \(-1\), same-unit cross-arm entries \(1\), and distinct cross-arm entries \(-1/3\).  Each row sum is therefore
\[
 3+3(-1)+1+3(-1/3)=0, \tag{15}
\]
which checks (5)--(14) over every assignment in the focal design.  At the opposite finite boundary, when \(n=4\) and \(G=K_4\), one has \(F_G(A)=\varnothing\) for every singleton \(A\); all trials fail and every \(\pi_u=0\).  This verifies that the positivity premise is indispensable rather than a removable boundary convention.
\end{proof}
\par\smallskip\noindent \textit{Justification.} This finite-design lemma separates the capsule identities that are exact from the unresolved random-regular operator upper bound. \textit{Gap.} Positive semidefiniteness, diagonal scale, incompatible-pair entries, and constant-mode cancellation do not control the heterogeneous compatible-pair block after reservoir acceptance conditioning. \textit{Consumer.} The capsule-kernel conjecture records these identities as its proved presolve capsule.

\begin{lemma}[lem:published-conflict-graph-inclusion-formulas]\label{lem:published-conflict-graph-inclusion-formulas}
Let \(B\) be a symmetric zero--one conflict matrix on \(V\) with diagonal entries one, let \(\Lambda=\|B\|_{\mathrm{op}}\), and let \(\sigma\) be an importance ordering.  Write \(M_\sigma(i)\) for the conflict neighbors of \(i\) preceding it in \(\sigma\).  In the Conflict Graph Design with sampling parameter \(\rho>1\), let \(E_{i,a}\) denote the event that unit \(i\) requests endpoint arm \(a\) and every member of \(M_\sigma(i)\) requests the null option.  Then
\[
\Pr(E_{i,a})=\frac{1}{2\rho\Lambda}
\left(1-\frac{1}{\rho\Lambda}\right)^{|M_\sigma(i)|}.
\]
For distinct \(i,j\), \(\Pr(E_{i,a}\cap E_{j,b})=0\) when \(B_{ij}=1\), and otherwise
\[
\Pr(E_{i,a}\cap E_{j,b})=\frac{1}{(2\rho\Lambda)^2}
\left(1-\frac{1}{\rho\Lambda}\right)^{|M_\sigma(i)\cup M_\sigma(j)|}.
\]
\end{lemma}

\begin{lemma}[lem:published-conflict-graph-local-covariance]\label{lem:published-conflict-graph-local-covariance}
Under the Conflict Graph Design with sampling parameter two, let \(B\), \(\Lambda\), \(\sigma\), and \(E_{i,a}\) be as in the design's conflict-matrix construction.  Then
\[
\operatorname{Var}\!\left(\frac{\mathbf 1\{E_{i,a}\}}{\Pr(E_{i,a})}\right)\le5\Lambda.
\]
For distinct \(i,j\), the covariance of the two normalized desired-exposure indicators is \(-1\) when their conflict-graph distance is one and is zero when that distance is at least three.  When their distance is two, its absolute value is at most
\[
\frac{6}{\Lambda}\,
\bigl|\{\ell\in V:B_{i\ell}=B_{j\ell}=1\}\bigr|.
\]
\end{lemma}

\begin{lemma}[lem:published-conflict-graph-risk]\label{lem:published-conflict-graph-risk}
For the global endpoint contrast under arbitrary neighborhood interference, let \(B_G\) be the \(V\)-by-\(V\) matrix with \((B_G)_{ij}=\mathbf 1\{\operatorname{dist}_G(i,j)\le2\}\) and let \(\Lambda_G=\|B_G\|_{\mathrm{op}}\).  The published Conflict Graph Design is outcome-independent and its modified Horvitz--Thompson estimator, including the design's desired-exposure metadata, is unbiased for \(\tau_G\) and satisfies
\[
\operatorname{Var}(T_{\mathrm{CG}}\mid G)
\le \frac{12.5\Lambda_G}{n}
\left\{\frac1n\sum_{i\in V}y_{i1}^2+
\frac1n\sum_{i\in V}y_{i0}^2\right\}.
\]
\end{lemma}

\begin{lemma}[lem:gate-conflict-spectral-map]\label{lem:gate-conflict-spectral-map}
For every simple \(d\)-regular graph \(G\), define \(B_G\) by \((B_G)_{ij}=\mathbf 1\{\operatorname{dist}_G(i,j)\le2\}\).  This is the unit-level conflict matrix for the global endpoint contrast, while the adjacency matrix of \(H_G\) is
\[
\begin{pmatrix}0&B_G\\B_G&0\end{pmatrix}.
\]
Consequently,
\[
\lambda(H_G)=\|B_G\|_{\mathrm{op}}\le1+d^2.
\]
\end{lemma}
\begin{proof}
Fix \(i,j\in V\).  If \(\operatorname{dist}_G(i,j)\le2\), then the closed neighborhoods \(N_G[i]\) and \(N_G[j]\) intersect.  Requiring the treated endpoint at \(i\) forces every unit in the first closed neighborhood to treatment, whereas requiring the control endpoint at \(j\) forces every unit in the second closed neighborhood to control.  The two requirements are therefore incompatible.  If \(\operatorname{dist}_G(i,j)>2\), the two closed neighborhoods are disjoint, so any requested pair of endpoint arms can be realized by assigning the two neighborhoods separately.  Thus the unit-level global-endpoint conflict relation is exactly \((B_G)_{ij}=1\).

By \(\text{def:gate-conflict-graph}\), edges of \(H_G\) occur only between opposite-arm copies and follow precisely this relation.  Its adjacency matrix therefore has the displayed block form.  Since \(B_G\) is real symmetric, an eigenpair \(B_Gx=\mu x\) produces eigenpairs
\[
\begin{pmatrix}0&B_G\\B_G&0\end{pmatrix}
\binom{x}{x}=\mu\binom{x}{x},
\qquad
\begin{pmatrix}0&B_G\\B_G&0\end{pmatrix}
\binom{x}{-x}=-\mu\binom{x}{-x}.
\]
Hence the spectral radius of the block matrix is \(\|B_G\|_{\mathrm{op}}\), proving the equality with \(\lambda(H_G)\).

Finally, from a fixed vertex there are at most \(1\) vertices at distance zero, \(d\) at distance one, and \(d(d-1)\) at distance two.  Every row sum of \(B_G\) is thus at most
\[
1+d+d(d-1)=1+d^2. \tag{8}
\]
For a symmetric matrix, the operator norm is at most the maximum absolute row sum; applying this elementary induced-norm inequality to the nonnegative matrix \(B_G\) and using (8) proves \(\|B_G\|_{\mathrm{op}}\le1+d^2\).  As finite boundary checks, \(d=0\) gives \(B_G=I_n\) and equality \(\|B_G\|_{\mathrm{op}}=1\), while \(G=K_4\) gives \(B_G\) equal to the all-ones matrix and \(\|B_G\|_{\mathrm{op}}=4\le1+3^2\).
\end{proof}
\par\smallskip\noindent \textit{Justification.} This lemma identifies the spectral quantity in the published unit-level GATE design with the frozen two-copy comparator and bounds it deterministically on regular graphs. \textit{Gap.} The mapping supplies the published \(O(d^2/n)\) upper rate but contains no logarithmic covariance improvement. \textit{Consumer.} The published-design upper theorem uses it to convert Theorem 4.1 into a degree-explicit minimax bound.

\begin{theorem}[thm:published-order-bracket]\label{thm:published-order-bracket}
There exist universal constants \(0<c<C<\infty\) such that, for every feasible degree sequence satisfying \(d=o(\sqrt n)\), \[\Pr_{G_n}\!\left\{c\frac{1+d}{n}\le R_2(G)\le C\frac{1+d^2}{n}\right\}\longrightarrow1.\] For each fixed feasible \(d\), every revealed simple \(d\)-regular graph satisfies \(R_2(G)=\Theta_d(n^{-1})\). The displayed bracket does not identify an exact degree-dependent frontier, and neither a lower bound of order \(r(n,d)\) nor attainment by the capsule is asserted.
\end{theorem}
\begin{proof}
By the repaired theorem \(\text{thm:all-procedure-lower}\), there are a universal \(c>0\) and events \(L_n\) such that
\[
\Pr_{G_n}(L_n)\longrightarrow1,
\qquad
R_2(G)\ge c\frac{1+d}{n}\quad\text{on }L_n. \tag{20}
\]

We derive an upper bound that is valid for every revealed simple \(d\)-regular graph, including bounded \(d\).  Fix such a graph and a schedule \(y\in\mathcal M_2(G)\).  The cited finite-sample Conflict Graph Design result \(\text{lem:published-conflict-graph-risk}\) supplies an outcome-independent design and a modified Horvitz--Thompson estimator \(T_{\mathrm{CG}}\), possibly using the design's internal desired-exposure metadata, such that
\[
\mathbb E(T_{\mathrm{CG}}\mid G)=\tau_G(y) \tag{21}
\]
and
\[
\operatorname{Var}(T_{\mathrm{CG}}\mid G)
\le \frac{12.5\|B_G\|_{\mathrm{op}}}{n}
\left\{\frac1n\sum_{i\in V}y_{i1}^2+
\frac1n\sum_{i\in V}y_{i0}^2\right\}, \tag{22}
\]
where \((B_G)_{ij}=\mathbf1\{\operatorname{dist}_G(i,j)\le2\}\).

To place this estimator in the observed-data class of \(\text{def:minimax-risk}\), define
\[
T_{\mathrm{CG}}^{\mathrm{RB}}
=\mathbb E[T_{\mathrm{CG}}\mid G,Z,y(Z)]. \tag{23}
\]
All relevant assignment and auxiliary spaces are finite, and the auxiliary law is fixed by the observed graph.  Thus (23) is the same finite conditional sum for every schedule and is a measurable function of \((G,Z,y(Z))\), with any value assigned on assignments of design probability zero.  The tower property and (21) give
\[
\mathbb E(T_{\mathrm{CG}}^{\mathrm{RB}}\mid G)=\tau_G(y). \tag{24}
\]
The conditional variance identity gives
\[
\operatorname{Var}(T_{\mathrm{CG}}\mid G)
=\operatorname{Var}(T_{\mathrm{CG}}^{\mathrm{RB}}\mid G)
+\mathbb E[\operatorname{Var}(T_{\mathrm{CG}}\mid G,Z,y(Z))\mid G], \tag{25}
\]
so (22), (24), and the nonnegativity of the final term imply
\[
\mathbb E[(T_{\mathrm{CG}}^{\mathrm{RB}}-\tau_G)^2\mid G]
\le \frac{12.5\|B_G\|_{\mathrm{op}}}{n}
\left\{\frac1n\sum_i y_{i1}^2+
\frac1n\sum_i y_{i0}^2\right\}. \tag{26}
\]
By \(\text{def:endpoint-class}\), each average in braces is at most one.  The proved spectral map \(\text{lem:gate-conflict-spectral-map}\) gives \(\|B_G\|_{\mathrm{op}}=\lambda(H_G)\le1+d^2\).  Hence (26) yields, uniformly over the class,
\[
\sup_{y\in\mathcal M_2(G)}
\mathbb E[(T_{\mathrm{CG}}^{\mathrm{RB}}-\tau_G)^2\mid G]
\le25\frac{1+d^2}{n}. \tag{27}
\]
The design and estimator in (23) are admissible in \(\text{def:minimax-risk}\); evaluating its infimum at this pair proves the deterministic inequality
\[
R_2(G)\le25\frac{1+d^2}{n}. \tag{28}
\]

Intersecting the sure event (28) with \(L_n\), equations (20) and (28) prove the probability statement with \(C=25\).  If \(d\) is fixed, the deterministic lower branch of \(\text{lem:balanced-hole-bayes-lower}\) and (28) give, for every revealed graph,
\[
\frac1{16n}\le R_2(G)\le25\frac{1+d^2}{n}, \tag{29}
\]
which is exactly \(R_2(G)=\Theta_d(n^{-1})\).

For the smallest admissible case \(n=4,d=0\), (29) reads \(1/64\le R_2(G)\le25/4\) when the sharper value \(b(G)=2\) from the finite audit is used on the left; hence no positivity, division, or boundary case is lost at the no-interference elbow.  For growing \(d\), the ratio of the deterministic upper endpoint in (28) to the proved random-regular lower endpoint in (20) is of order at most \(d\), not of logarithmic order.  Obtaining the formerly claimed logarithmic-width bracket would require the unproved balanced-hole estimate (19), so neither that estimate nor capsule attainment is inferred from (20)--(29).
\end{proof}
\par\smallskip\noindent \textit{Justification.} This is the strongest two-sided conclusion supported by the exact published-class lower theorem and the closed published-design upper theorem. \textit{Gap.} For growing \(d\), the proved endpoints differ by a factor of order at most \(d\); neither endpoint is asserted to identify the exact minimax degree dependence. \textit{Consumer.} It is the paper's honest design-based minimax headline.

\begin{lemma}[lem:balanced-hole-bayes-lower]\label{lem:balanced-hole-bayes-lower}
Fix \(n\ge4\) and a revealed simple graph \(G\) on \(V\).  For \(z\in\{0,1\}^V\) and \(a\in\{0,1\}\), define
\[
E_a(z)=\{i\in V:z_j=a\text{ for every }j\in N_G[i]\},
\qquad
b(G)=\max_{z\in\{0,1\}^V}\min\{|E_1(z)|,|E_0(z)|\}.
\]
Then the unrestricted design-based minimax risk satisfies the deterministic inequalities
\[
R_2(G)\ge \frac{1}{32\max\{1,b(G)\}}\ge \frac{1}{16n}.
\]
The first inequality holds simultaneously against every outcome-independent assignment law and every measurable, possibly biased or randomized estimator on \(\mathcal M_2(G)\).
\end{lemma}
\begin{proof}
Put \(B=\max\{1,b(G)\}\).  We first record two deterministic consequences of the definitions.  Since \(i\in N_G[i]\), a unit cannot belong to both \(E_1(z)\) and \(E_0(z)\).  Hence, for every assignment \(z\),
\[
\min\{|E_1(z)|,|E_0(z)|\}
\le \frac{|E_1(z)|+|E_0(z)|}{2}
\le \frac n2. \tag{1}
\]
Taking the maximum over \(z\) gives \(b(G)\le n/2\).  Because \(n\ge4\), it follows also that
\[
1\le B\le \frac n2. \tag{2}
\]
Moreover, if \(i\in E_1(z)\) and \(j\in E_0(z)\), then \(N_G[i]\cap N_G[j]=\varnothing\): a vertex in the intersection would be required by the two exposure definitions to have both assignment values.  If \(\operatorname{dist}_G(i,j)\le2\), the two closed neighborhoods intersect, at \(i\), at \(j\), or at an intermediate vertex according as the distance is zero, one, or two.  Consequently
\[
i\in E_1(z),\ j\in E_0(z)
\quad\Longrightarrow\quad
\operatorname{dist}_G(i,j)>2. \tag{2a}
\]
Thus \(b(G)\) is indeed the largest balanced two-sided radius-two hole generated by an assignment.

We prove the first inequality by a Bayes reduction whose prior is common to all designs and estimators.  For each \(a\in\{0,1\}\), independently draw a Rademacher sign \(B_a\).  Conditional on \(B_a\), independently over \(i\in V\) draw \(\xi_{ia}\in\{-1,1\}\) with
\[
\Pr(\xi_{ia}=1\mid B_a)=\frac{1+\delta B_a}{2},
\qquad
\delta^2=\frac{1}{4B}. \tag{3}
\]
Equation (2) implies \(0<\delta\le1/2\), so the probabilities in (3) are strictly between zero and one.  Define a random fixed potential-outcome schedule by
\[
y_i(z)=
\begin{cases}
\xi_{i1},&z_j=1\text{ for every }j\in N_G[i],\\
\xi_{i0},&z_j=0\text{ for every }j\in N_G[i],\\
0,&\text{otherwise}.
\end{cases} \tag{4}
\]
The first two cases are disjoint because \(N_G[i]\) contains \(i\).  Formula (4) depends on \(z\) only through \(z|_{N_G[i]}\), so it satisfies \(\text{ass:ani}\).  Moreover,
\[
y_{i1}=\xi_{i1},\qquad y_{i0}=\xi_{i0},
\qquad
\frac1n\sum_{i\in V}y_{ia}^2=1\quad(a=0,1). \tag{5}
\]
Thus (5), \(\text{ass:treated-endpoint-l2}\), \(\text{ass:control-endpoint-l2}\), and \(\text{def:endpoint-class}\) show that every realization of the prior belongs to \(\mathcal M_2(G)\), including the boundary of both moment balls.

Fix an arbitrary outcome-independent assignment law and condition on \(Z=z\).  By (4), the observed outcomes reveal precisely \((\xi_{i1}:i\in E_1(z))\), \((\xi_{i0}:i\in E_0(z))\), and deterministic zeros at all remaining units.  Select an arm \(a_*(z)\) attaining the smaller exposure count and put
\[
m=|E_{a_*(z)}(z)|\le b(G)\le B,
\qquad L=n-m. \tag{6}
\]
We next lower-bound the posterior uncertainty in this arm.

Let \(P_+^{(m)}\) and \(P_-^{(m)}\) denote the laws of its \(m\) revealed signs conditional on \(B_{a_*}=1\) and \(B_{a_*}=-1\), respectively.  For one sign, direct summation gives Hellinger affinity \(\sqrt{1-\delta^2}\); independence conditional on \(B_{a_*}\) therefore gives
\[
\rho_m=\sum_x\sqrt{P_+^{(m)}(x)P_-^{(m)}(x)}
=(1-\delta^2)^{m/2}. \tag{7}
\]
Writing \(p_+\) and \(p_-\) for the two mass functions, Cauchy--Schwarz yields
\[
\begin{aligned}
\operatorname{TV}(P_+^{(m)},P_-^{(m)})
&=\frac12\sum_x|\sqrt{p_+(x)}-\sqrt{p_-(x)}|
                    (\sqrt{p_+(x)}+\sqrt{p_-(x)})\\
&\le \sqrt{1-\rho_m^2}
=\sqrt{1-(1-\delta^2)^m}
\le \sqrt{m\delta^2}
\le\frac12. \tag{8}
\end{aligned}
\]
The penultimate inequality follows by expanding \(1-(1-x)^m\le mx\) for \(x\in[0,1]\), and the last follows from (3) and (6).

Let \(q(x)=\mathbb E[B_{a_*}\mid x]\) under the equal mixture of \(P_+^{(m)}\) and \(P_-^{(m)}\).  Bayes' rule gives
\[
q(x)=\frac{p_+(x)-p_-(x)}{p_+(x)+p_-(x)}
\]
whenever the denominator is positive.  Integrating against the mixture mass \((p_++p_-)/2\) gives
\[
\mathbb E|q|=\operatorname{TV}(P_+^{(m)},P_-^{(m)}). \tag{9}
\]
Since \(|q|\le1\), equations (8)--(9) imply
\[
\mathbb E\!\left[\operatorname{Var}(B_{a_*}\mid x)\right]
=1-\mathbb E[q(x)^2]
\ge1-\mathbb E|q(x)|
\ge\frac12. \tag{10}
\]

Given \(B_{a_*}\), the \(L\) unobserved endpoint signs in this arm remain conditionally independent, each with mean \(\delta B_{a_*}\) and variance \(1-\delta^2\).  Applying the conditional variance identity first given the revealed signs and then given \(B_{a_*}\), and using (10), gives
\[
\begin{aligned}
&\mathbb E\!\left[
\operatorname{Var}\!\left(
\frac1n\sum_{i\in V}\xi_{i,a_*}\,\middle|\,
z,\text{ observed outcomes}
\right)\,\middle|\,z\right]\\
&\quad=\frac1{n^2}\mathbb E\!\left[
L(1-\delta^2)+L^2\delta^2
\operatorname{Var}(B_{a_*}\mid x)\,\middle|\,z\right]\\
&\quad\ge\frac{L^2\delta^2}{2n^2}. \tag{11}
\end{aligned}
\]
By (1), (2), and (6), \(L=n-m\ge n-b(G)\ge n/2\).  Substitution of (3) into (11) therefore yields
\[
\mathbb E\!\left[
\operatorname{Var}\!\left(
\frac1n\sum_i\xi_{i,a_*}\,\middle|\,
z,\text{ observed outcomes}
\right)\,\middle|\,z\right]
\ge \frac1{32B}. \tag{12}
\]

The two arm priors are independent, and conditional on \(z\) their likelihood factors by arm.  Consequently the posterior covariance between the two endpoint means is zero.  Since
\[
\tau_G(y)=\frac1n\sum_i(\xi_{i1}-\xi_{i0}), \tag{13}
\]
the posterior variance of \(\tau_G(y)\) is the sum of the two arm posterior variances; its expectation conditional on \(z\) is therefore at least the lower bound in (12).

Let \(\mathcal F\) contain \(G\), \(Z\), \(y(Z)\), and any estimator random seed.  Outcome independence of the design and independence of the estimator seed from the prior ensure that neither carries additional information about the endpoint signs beyond the observations just analyzed.  For every \(\mathcal F\)-measurable estimator \(T_*\), completing the square conditionally gives
\[
\mathbb E[(T_*-\tau_G)^2]
=\mathbb E[(T_*-\mathbb E[\tau_G\mid\mathcal F])^2]
+\mathbb E[\operatorname{Var}(\tau_G\mid\mathcal F)]
\ge\frac1{32B}. \tag{14}
\]
The Bayes risk in (14) is no larger than the supremum risk over the support of the prior, which is contained in \(\mathcal M_2(G)\) by (5).  Since the same prior was used for every admissible design--estimator pair, taking the infimum specified by \(\text{def:minimax-risk}\) proves
\[
R_2(G)\ge\frac1{32B}
=\frac1{32\max\{1,b(G)\}}. \tag{15}
\]
Finally, (2) turns (15) into \(R_2(G)\ge1/(16n)\).

For the required smallest-instance audit, take \(n=4\) and \(d=0\).  Then \(E_1(z)=\{i:z_i=1\}\) and \(E_0(z)=\{i:z_i=0\}\), so enumeration by the treated count gives \(b(G)=2\).  Formula (15) gives \(R_2(G)\ge1/64\).  At the opposite finite boundary \(G=K_4\), every nonconstant assignment has \(E_1(z)=E_0(z)=\varnothing\), while a constant assignment exposes only one arm; hence \(b(G)=0\) and (15) gives \(R_2(G)\ge1/32\).  Both checks agree with (1)--(15), including the \(\max\{1,b(G)\}\) convention at \(b(G)=0\).
\end{proof}
\par\smallskip\noindent \textit{Justification.} This deterministic Bayes certificate cleanly separates the valid all-procedure reduction from the unresolved random-regular balanced-hole estimate. \textit{Gap.} The lemma does not assert a random-regular upper bound on \(b(G)\); that is the missing combinatorial step behind the discarded logarithmic-scale lower edge. \textit{Consumer.} The repaired lower theorem uses its universal \(\Omega(n^{-1})\) branch, and the repaired bracket uses it at the bounded-degree elbow.

\begin{lemma}[lem:published-random-regular-gate-lower]\label{lem:published-random-regular-gate-lower}
For the global average treatment effect under arbitrary neighborhood interference and separate empirical second-moment bounds on its two endpoint arms, if \(G_n\) is uniformly random among simple labeled \(d\)-regular graphs and \(d=o(\sqrt n)\), then there is a universal constant \(c_0>0\) such that
\[
\Pr_{G_n}\!\left\{R_2(G)\ge c_0\frac d n\right\}\longrightarrow1.
\]
\end{lemma}
\par\smallskip\noindent \textit{Justification.} This is the matching primary-source lower bound that remains after the unsupported stronger balanced-hole switching claim is removed. \textit{Gap.} The cited corollary proves the \(d/n\) lower edge for random-regular GATE; it does not prove the stronger \((d+1)^2/[n\log(d+2)]\) edge. \textit{Consumer.} The repaired all-procedure lower theorem combines this growing-degree edge with the deterministic \(1/n\) branch.

\section{Interpretation}

The supported bracket rules out risk below the published lower order and gives the ordinary \(n^{-1}\) benchmark at every fixed degree, but it does not determine the growing-degree frontier. A valid balanced-hole theorem would strengthen the all-procedure lower edge to \(r(n,d)\). Independently, an \(O(n/k)\) capsule norm certificate would yield a sharper observed-data upper construction. Both advances remain separate obligations.

\section*{Technical internal limitation (diagnostic only)}

Two nontrivial steps remain unresolved. First, the proposed random-regular balanced-hole estimate needs uniform switching counts that explicitly control simple-graph invalid moves, exceptional internal-edge strata, and the available half-edge event. Second, the capsule's accepted compatible-pair covariance block contains the random denominator \(\binom{|F_G(A)|}{k}\), and neither entrywise bounds nor pointwise Monte Carlo supplies the required operator-norm certificate or efficient conditional Rao--Blackwell weights. No switching, covariance, overlap, or computational premise is added to assume either crux.

\section{Honest scope}

The paper does not prove the strengthened random-regular lower edge, an exact minimax frontier, or capsule attainment at \(d^2/[n\log d]\). It does prove or retain the supported published-order bracket and the fixed-degree \(\Theta_d(n^{-1})\) benchmark. The capsule assignment law is total, but its normalized kernel and estimator audit are asserted only after finite enumeration verifies every \(\pi_u>0\); \(G=K_4\) at \(n=4\) is an explicit zero-inclusion witness. Exact enumeration may be exponential, and scalable certification remains open. No studentized central limit theorem, bounded-outcome model, additive interference model, or real-network random-regular approximation is claimed.

\begin{thebibliography}{99}

  \bibitem{KandirosHarshawSavje2026DesignBasedMinimax} Vardis Kandiros, Christopher Harshaw, and Fredrik Savje. A Design-Based Minimax Theory for Network Experiments. arXiv:2608.04909, 2026.

  \bibitem{KandirosPipisDaskalakisHarshaw2024ConflictGraph} Vardis Kandiros, Charilaos Pipis, Constantinos Daskalakis, and Christopher Harshaw. The Conflict Graph Design: Estimating Causal Effects under Arbitrary Neighborhood Interference. arXiv:2411.10908, 2024.

  \bibitem{LiLi2026PseudoRandomDesign} Shuangning Li and Shuangping Li. Minimax-Optimal Experimental Design for Network Interference on Pseudo-Random Graphs. IMSI Recent Advances in Random Networks slides, 2026.

  \bibitem{HarshawSavjeWang2025Riesz} Christopher Harshaw, Fredrik S\"avje, and Yitan Wang. A General Design-Based Framework and Estimator for Randomized Experiments. arXiv:2210.08698v3, 14 August 2025.

  \bibitem{CaiZhangAiroldi2024IndependentSet} Chencheng Cai, Xu Zhang, and Edoardo M. Airoldi. Independent-Set Design of Experiments for Estimating Treatment and Spillover Effects under Network Interference. International Conference on Learning Representations, 2024.

  \bibitem{PerkinsWang2024BalancedIndependentSets} Will Perkins and Yuzhou Wang. On the Hardness of Finding Balanced Independent Sets in Random Bipartite Graphs. Proceedings of the 2024 ACM--SIAM Symposium on Discrete Algorithms, pages 2376--2397, 2024.

  \bibitem{NarangPerkinsWangWee2026FixedSlice} Ijay Narang, Will Perkins, Yuzhou Wang, and Timothy L. H. Wee. Computational Thresholds for Balanced and Fixed-Slice Independent Sets in Bipartite Graphs. arXiv:2608.02503, 2026.

  \bibitem{AronowSamii2017GeneralInterference} Peter M. Aronow and Cyrus Samii. Estimating Average Causal Effects under General Interference, with Application to a Social Network Experiment. Annals of Applied Statistics 11(4):1912--1947, 2017.

  \bibitem{UganderYin2023RGCR} Johan Ugander and Hao Yin. Randomized Graph Cluster Randomization. Journal of Causal Inference 11(1), 2023.

  \bibitem{KarrerEtAl2021NetworkScale} Brian Karrer, Liang Shi, Mehrdad Bhole, Matt Goldman, Tim Palmer, Chris Gelman, Murat Konutgan, and Feng Sun. Network Experimentation at Scale. Proceedings of the 27th ACM SIGKDD Conference on Knowledge Discovery and Data Mining, 2021.

  \bibitem{Pippenger1989RSA} Nicholas Pippenger. Random Sequential Adsorption on Graphs. SIAM Journal on Discrete Mathematics 2:393--401, 1989.

  \bibitem{KrivelevichMesarosMichaeliShikhelman2024Greedy} Michael Krivelevich, Tamas Meszaros, Peleg Michaeli, and Clara Shikhelman. Greedy Maximal Independent Sets via Local Limits. Random Structures and Algorithms 64(2):433--470, 2024.

  \bibitem{Kallus2021OptimalRandomization} Nathan Kallus. On the Optimality of Randomization in Experimental Design: How to Randomize for Minimax Variance and Design-Based Inference. Journal of the Royal Statistical Society, Series B 83(2):404--409, 2021.

\end{thebibliography}

\end{document}